% paralleles.tex — f163r, vv. 43–79 : loci paralleli. Généré par .travail/f163r_lp_redac/rediger.py
% à partir des collectes vérifiées (paralleles_43-61.json, paralleles_62-79.json) et des parallèles
% d'Austin–Olson 2004 retrouvés dans les sources ; règles : edition/CONVENTIONS.md § 4
\lp{43–45 ἐχέτω δὲ πνοὰς … γλαυκόν]}{\textit{Test.} Suid.~ε~3997 ἐχέτω· κρατείτω. Ἀριστοφάνης· «ἐχέτω δὲ πνοὰς νήνεμος αἰθήρ, κῦμα δὲ πόντου μὴ κελαδείτω γλαυκόν.» ; \textit{idem} Zonar. col.~941~Tittmann}
\lp{43–48 ἐχέτω δὲ πνοὰς … μὴ λυέσθων]}{Ar.~\textit{Av.}~777–778 πτῆξε δὲ φῦλά τε ποικίλα θηρῶν, κύματά τ᾽ ἔσβεσε νήνεμος αἴθρη ; Eur.~\textit{Ba.}~1084–1085 σίγησε δ᾽ αἰθήρ, σῖγα δ᾽ ὕλιμος νάπη φύλλ᾽ εἶχε, θηρῶν δ᾽ οὐκ ἂν ἤκουσας βοήν ; Eur.~\textit{Hel.}~1455–1458 ὅταν αὔραις πέλαγος ἀνήνεμον ᾖ, γλαυκὰ δὲ Πόντου θυγάτηρ Γαλάνεια τάδ᾽ εἴπῃ ; Eur.~\textit{IA}~9–11 οὔκουν φθόγγος γ᾽ οὔτ᾽ ὀρνίθων οὔτε θαλάσσης· σιγαὶ δ᾽ ἀνέμων τόνδε κατ᾽ Εὔριπον ἔχουσιν ; Aesch.~\textit{Ag.}~565–566 εὖτε πόντος ἐν μεσημβριναῖς κοίταις ἀκύμων νηνέμοις εὕδοι πεσών ; \textit{Od.}~12.168–169 αὐτίκ᾽ ἔπειτ᾽ ἄνεμος μὲν ἐπαύσατο ἠδὲ γαλήνη ἔπλετο νηνεμίη, κοίμησε δὲ κύματα δαίμων ; Alcm.~fr.~60~Bergk (=~89~PMG) Εὕδουσιν δ᾽ ὀρέων κορυφαί τε καὶ φάραγγες … θῆρές τ᾽ ὀρεσκῷοι καὶ γένος μελισσᾶν καὶ κνώδαλ᾽ ἐν βένθεσι πορφυρέας ἁλός· εὕδουσιν δ᾽ οἰωνῶν φῦλα τανυπτερύγων ; Simon.~fr.~37~Bergk ap. D.H.~\textit{Comp.}~26 κέλομαι· εὗδε βρέφος, εὑδέτω δὲ πόντος, εὑδέτω ἄμετρον κακόν ; Call.~\textit{Ap.}~17–18 εὐφημεῖτ᾽ ἀίοντες ἐπ᾽ Ἀπόλλωνος ἀοιδῆι. εὐφημεῖ καὶ πόντος, ὅτε κλείουσιν ἀοιδοί ; Theoc.~2.38 ἠνίδε σιγῇ μὲν πόντος, σιγῶντι δ᾽ ἀῆται ; Limen.~\textit{Paean.} (Jan, Mus. Script. Gr. Suppl. p.~24, ped.~40–51) νηνέμους δ᾽ ἔσχεν αἰθὴρ ἀελλωῶν ταχυπετεῖς δρόμους, λῆξε δὲ βαρύβρομον Νηηρέως ζαμενὲς οεῖδμ᾽ ἠδὲ μέγας Ὠκεανός ; Mesomed.~\textit{Hymn.Sol.}~1–6 (Jan, Mus. Script. Gr. Suppl. p.~48–49) Εὐφαμείτω πᾶς αἰθήρ, γῆ καὶ πόντος καὶ πνοαί, οὔρεα τέμπεα σιγάτω, ἦχοι φθόγγοι τ᾽ ὀρνίθων. μέλλει γὰρ πρὸς ἡμᾶς βαίνειν Φοῖβος ἀκερσεκόμας εὐχαίτας ; Luc.~\textit{Trag.}~129–133 Σῖγα μὲν αἰθὴρ νήνεμος ἔστω, καὶ πᾶς ποδαγρῶν εὐφημείτω. ἴδε, πρὸς θυμέλας ἡ κλινοχαρὴς βαίνει δαίμων ; Pl.~\textit{Smp.}~197c οὗτός ἐστιν ὁ ποιῶν εἰρήνην μὲν ἐν ἀνθρώποις, πελάγει δὲ γαλήνην νηνεμίαν, ἀνέμων κοίτην ὕπνον τ᾽ ἐνὶ κήδει ; Soph.~\textit{Aj.}~674–675 δεινῶν τ᾽ ἄημα πνευμάτων ἐκοίμισε στένοντα πόντον ; A.R.~3.749–750 οὐ θρόος ἦεν ἠχήεις· σιγὴ δὲ μελαινομένην ἔχεν ὄρφνην ; Synes.~\textit{Hymn.}~3.72–83 (PG~66, col.~1594–1595) Εὐφαμείτω / Αἰθὴρ καὶ γᾶ· / Στάτω πόντος, / Στάτω δ᾽ ἀήρ· / Λήγετε, πνοαὶ / Βαλίων ἀνέμων, / Λήγετε, ῥιπαὶ / Γυρῶν ῥοθίων, / Ποταμῶν προχοαί, / Κρανααὶ λιβάδες. / Ἐχέτω σιγὰ / Κόσμου λαγόνας ; Verg.~\textit{Aen.}~4.522–527 Nox erat, et placidum carpebant fessa soporem corpora per terras, silvaeque et saeva quierant aequora: … cum tacet omnis ager, pecudes pictaeque volucres ; Verg.~\textit{Ecl.}~9.57–58 et nunc omne tibi stratum silet aequor, et omnes, aspice, ventosi ceciderunt murmuris aurae ; Hdt.~7.193.1 ὡς ἐπαύσατό τε ὁ ἄνεμος καὶ τὸ κῦμα ἔστρωτο ; Philostr.~\textit{Im.}~2.15.1 θέλγει τὴν θάλατταν Ὀρφεὺς ᾅδων, ἡ δὲ ἀκούει καὶ ὑπὸ τῇ ᾠδῇ κεῖται ὁ Πόντος}
\lp{43 ἐχέτω δὲ πνοάς]}{Ar.~\textit{Av.}~1397–1398 ἅμ᾽ ἀνέμων πνοαῖσι βαίην. — νὴ τὸν Δί᾽ ἦ ᾽γώ σου καταπαύσω τὰς πνοάς ; Σ Ar.~\textit{Thesm.}~43 ἐπεχέτω ; Eur.~\textit{Hec.}~900–901 νῦν δ᾽, οὐ γὰρ ἵησ᾽ οὐρίους πνοὰς θεός, μένειν ἀνάγκη ; \textit{Od.}~5.383–384 ἦ τοι τῶν ἄλλων ἀνέμων κατέδησε κελεύθους, παύσασθαι δ᾽ ἐκέλευσε καὶ εὐνηθῆναι ἅπαντας ; Eur.~\textit{Or.}~700 ἴσως ἂν ἐκπνεύσειεν· ἢν δ᾽ ἀνῇ πνοάς ; Zen.~3.17 Δικτύῳ ἄνεμον θηρᾷς; ἐπὶ τῶν μάτην καὶ ἀνοήτως τὸ ποιούντων}
\lp{43 πνοάς]}{Eur.~\textit{Phoen.}~211–213 Ζεφύρου πνοαῖς ἱππεύσαντος, ἐν οὐρανῷ κάλλιστον κελάδημα ; Hsch.~α~1518 ἀήτας· πνοάς (Ξ 254) ; Eur.~\textit{Andr.}~479 πνοαὶ δ᾽ ὅταν φέρωσι ναυτίλους θοαί ; Aesch.~\textit{Ag.}~192 πνοαὶ δ᾽ ἀπὸ Στρυμόνος μολοῦσαι ; Soph.~\textit{El.}~435 ἀλλ᾽ ἢ πνοαῖσιν ἢ βαθυσκαφεῖ κόνει ; \textit{Il.}~11.622 στάντε ποτὶ πνοιὴν παρὰ θῖν᾽ ἁλός ; Pi.~\textit{O.}~3.31–32 κείναν χθόνα πνοιᾶς ὄπιθεν Βορέα ψυχροῦ ; Hsch.~α~8336 αὖραι· πνοαί ; Hsch.~π~2649 πνοή· ἀνάψυξις}
\lp{43 νήνεμος αἰθήρ]}{\textit{Il.}~8.555–556 ἄστρα φαεινὴν ἀμφὶ σελήνην φαίνετ᾽ ἀριπρεπέα, ὅτε τ᾽ ἔπλετο νήνεμος αἰθήρ ; Theoc.~7.57–58 χἁλκυόνες στορεσεῦντι τὰ κύματα τάν τε θάλασσαν τόν τε νότον τόν τ᾽ εὖρον, ὃς ἔσχατα φυκία κινεῖ ; A.R.~1.1154–1155 ἀμφὶ γὰρ αἰθὴρ νήνεμος ἐστόρεσεν δίνας, κατὰ δ᾽ εὔνασε πόντον ; Zen.~1.39 Αἰθέρα νήνεμον ἐρέσσειν; ἐπὶ τῶν μάτην πονούντων ; \textit{idem} Diogenian.~1.38 (αἱρήσεις), Suid.~αι~114 (ἐρέσσεις) ; Ar.~\textit{Ran.}~1003–1003a ἡνίκ᾽ ἂν τὸ πνεῦμα λεῖον καὶ καθεστηκὸς λάβῃς ; Eur.~\textit{Or.}~279 ἐκ κυμάτων γὰρ αὖθις αὖ γαλήν᾽ ὁρῶ ; \textit{idem} Ar.~\textit{Ran.}~304 (γαλῆν), Apostol.~6.88d (om. αὖ ; γαλήνην ; add. Εὐριπίδου) ; \textit{Il.}~5.522–524 ἀλλ᾽ ἔμενον νεφέλῃσιν ἐοικότες ἅς τε Κρονίων νηνεμίης ἔστησεν ἐπ᾽ ἀκροπόλοισιν ὄρεσσιν ἀτρέμας, ὄφρ᾽ εὕδῃσι μένος Βορέαο ; \textit{Od.}~5.391–392 καὶ τότ᾽ ἔπειτ᾽ ἄνεμος μὲν ἐπαύσατο ἠδὲ γαλήνη ἔπλετο νηνεμίη ; Hsch.~α~1885 αἴθρη· αἰθρία. αὐγή. εὐδία. ἀὴρ καθαρός (ζ 44 . .) ; Hsch.~α~1884 αἰθρία· εὐδία (Xen. Hell. 7, 1, 31 . .) ; Hsch.~γ~92 γαληνόν· ἥσυχον. ἱλαρόν. εὔδιον ; Hsch.~ε~716 εἱαμενόν· νήνεμον. κοῖλον. [βοτανώδη] ; Hsch.~η~916 ἡσυχία· γαλήνη, ἡρεμία (Eur. Alc. 77) ; Hsch.~σ~1933 στορεύς· γαληνοποιός ; Zen.~1.38 Αἰγιαλῷ λαλεῖς: ἐπὶ τῶν ἀνηκούστων. Ὁμοία δὲ αὕτη καὶ τῇ, Ἀνέμῳ διαλέγῃ ; Erasm.~\textit{Adag.}~1.4.85 Vento loqueris. Suidas addit \& hoc: Ἀνέμῳ διαλέγῃ, i. e. Vento loqueris. Pro eo quod est, inaniter}
\lp{43 νήνεμος αἰθήρ (ἀλκυονίδες ἡμέραι)]}{Σ Ar.~\textit{Ran.}~1309 ἀλκυόνες: Χαρακτηρίζει τὰ Εὐριπίδου μέλη ὡς ἐκλελυμένα. [ἀλκυόνες δὲ ἐπάνω τοῦ τῆς θαλάσσης ὕδατος καθήμεναι … καθ᾽ ὃν δὲ καιρὸν τοῦτο ποιοῦσιν, ἀλκυονίτιδες αὗται αἱ ἡμέραι καλοῦνται, καὶ ἔστιν ἐν αὐταῖς μεγίστη ἡσυχία ἀνέμων τε καὶ κυμάτων.] ; Ar.~\textit{Av.}~1594 ἀλκυονίδας τ᾽ ἂν ἤγεθ᾽ ἡμέρας ἀεί ; Simon.~fr.~12~Bergk ap. Arist. HA~5.8 (542b) ὡς ὁπόταν χειμέριον κατὰ μῆνα πινύσκῃ Ζεὺς ἤματα τεσσαρακαίδεκα, λαθάνεμόν τέ μιν ὥραν καλέουσιν ἐπιχθόνιοι, ἱερὰν παιδοτρόφον ποικίλας ἁλκυόνος ; Hsch.~α~3100 ἁλκυονίδες· ἡμέραι τινὲς τὸν ἀριθμὸν ιδʹ γαληναί, ἐν αἷς νοσσεύει ἡ ἁλκυών ; Suid.~α~1298 Ἀλκυονίδες ἡμέραι: αἱ εὐδιειναί. … αἱ δὲ νήνεμοι καὶ γαλήνην ἔχουσαι ἡμέραι Ἀλκυονίδες καλοῦνται. καὶ Ἀλκυόνειος ἡμέρα ; Erasm.~\textit{Adag.}~2.6.52 Halcedonia sunt apud forum. … Significatur hac voce tranquillitas \& silentium. … Ἀλκυονίδας ἡμέρας ἄγεις. i. e. Halcyonios agis dies. De tranquillam \& ociosam agentibus vitam. Aristophanes in Avibus: Ἀλκυονίδας τ᾽ ἦγεθ᾽ ἡμέρας ἀεί ; Plaut.~\textit{Poen.}~355–356 nisi illam mihi tam tranquillam facis quam mare olimst, quom ibi alcedo pullos educit suos ; Tosi n°~1702 Tam tranquillam.../ quam mare. Aussi calme que la mer}
\lp{43 νήνεμος]}{Eur.~\textit{Hec.}~532–533 Σιγᾶτ᾽, Ἀχαιοί, σῖγα πᾶς ἔστω λεώς, σίγα σιώπα· νήνεμον δ᾽ ἔστησ᾽ ὄχλον ; Σ Eur.~\textit{Hec.}~533 νήνεμον: ἀθόρυβον. … ἥσυχον. … ἀπὸ μεταφορᾶς τῶν ἀνέμων ; Hsch.~ν~488 νηνεμίης· πνοῆς ἐκτός. εὐδίας (Ε 523) ; Hsch.~ν~489 νήνεμον· ἄνευ ἀνέμου. εὔπνουν. εὔδιον. ἥσυχον. ἄπνουν. καὶ ἀνήνεμον τὸ αὐτό ; Eur.~\textit{IT}~1412 εἰ μὴ γὰρ οἶδμα νήνεμον γενήσεται ; Aesch.~\textit{Ag.}~738–739 φρόνημα μὲν νηνέμου γαλάνας ; A.R.~2.162 περὶ δέ σφιν ἰαίνετο νήνεμος ἀκτὴ ; A.R.~2.660–661 ὁμῶς δ᾽ ἐπὶ ἤματι νύκτα / νήνεμον ἀκαμάτῃσιν ἐπερρώοντ᾽ ἐλάτῃσιν ; Hsch.~α~1433 ἀζαλής· νήνεμος, γαληνός ; Hsch.~γ~287 γελαρής· γαληνή. Λάκωνες ; Hsch.~ι~1000 ἴσσος· γαλήνη ; Hsch.~ν~487 νηνεμία· γαλήνη, ἀνέμων ⟨ἀπουσία⟩ ; Hsch.~ν~490 νηνεμούμενον· καθεσταμένον ἐκ ταραχῆς. εὔδιον. ἥσυχον ; Phot.~ν~190 νήνεμος· ἄνεμον οὐκ ἔχων ; \textit{idem} Suid.~ν~319}
\lp{43 αἰθήρ]}{Ar.~\textit{Thesm.}~14 αἰθὴρ γὰρ ὅτε τὰ πρῶτα διεχωρίζετο ; Ar.~\textit{Thesm.}~272 ὄμνυμι τοίνυν αἰθέρ᾽ οἴκησιν Διός ; Ar.~\textit{Ran.}~892 αἰθὴρ ἐμὸν βόσκημα καὶ γλώσσης στρόφιγξ ; Hsch.~α~1860 αἰθήρ· ὁ ὑπὲρ τὰ νέφη τόπος}
\lp{44 κῦμα δὲ πόντου]}{Aesch.~\textit{PV}~1048–1049 κῦμα δὲ πόντου τραχεῖ ῥοθίῳ συγχώσειεν τῶν οὐρανίων ἄστρων διόδους ; Ar.~\textit{Ran.}~1309–1310 ἀλκυόνες, αἳ παρ᾽ ἀενάοις θαλάσσης κύμασι στωμύλλετε ; Ar.~\textit{Eq.}~432–433 ἐγὼ δὲ συστείλας γε τοὺς ἀλλᾶντας εἶτ᾽ ἀφήσω / κατὰ κῦμ᾽ ἐμαυτὸν οὔριον ; Ar.~\textit{Ran.}~704 τὴν πόλιν καὶ ταῦτ᾽ ἔχοντες κυμάτων ἐν ἀγκάλαις ; Eur.~\textit{Or.}~343–344 κατέκλυσεν δεινῶν πόνων ὡς πόντου λάβροις ὀλεθρίοισιν ἐν κύμασιν ; Hes.~\textit{Op.}~691 δεινὸν γὰρ πόντου μετὰ κύμασι πήματι κύρσαι ; Hsch.~κ~4523 κῦμα· ὕδατος ἐπιδρομὴ βιαία. καὶ τὸ ἔτι ὂν ἐν τῇ κοιλίᾳ. καὶ τὸ γεννώμενον κύημα ; Hsch.~κ~4534 κύματα· ὑδάτων ἐπιδρομαὶ βίαιοι. κυήματα. τὰ γεννώμενα. καὶ τὰ ἐν ταῖς κοιλίαις ἔτι ὄντα· καὶ δένδρα, καὶ ὅρη ; Hsch.~π~2996 πόντος· θάλασσα, πέλαγος (Β 210 . .)}
\lp{44 μὴ κελαδείτω]}{Ar.~\textit{Nub.}~283–284 καὶ ποταμῶν ζαθέων κελαδήματα, καὶ πόντον κελάδοντα βαρύβρομον ; Hsch.~κ~4529 κύματι κωφῷ· ἀήχῳ (Ξ 16) ; Hsch.~κ~4531 κύματι κωφῷ· τῷ μὴ ἠχοῦντι, ἀλλ᾽ ἀρχομένῳ μεγεθοῦσθαι (Ξ 16) ; \textit{Od.}~2.420–421 ἀκραῆ Ζέφυρον, κελάδοντ᾽ ἐπὶ οἴνοπα πόντον ; Orac. ap. Paus.~10.37.6 πρίν κεν ἐμῷ τεμένει κυανώπιδος Ἀμφιτρίτης κῦμα ποτικλύζῃ κελαδοῦν ἐπὶ οἴνοπα πόντον ; Hsch.~α~2432 ἀκλυδώνιστον· ἀτάραχον ; Hsch.~α~2691 ἀκύμονα· ἀτάραχον (Cyr. in Ioh. 1, 343, 15) ; Hsch.~κ~2127 κελαδεινῆς· θορυβώδους (Π 183) ; Hsch.~κ~2128 κελάδησαν· ἤχησαν (Θ 542) ; Hsch.~κ~2129 κέλαδος· ἦχος μουσικός, κραυγή, βοή, θόρυβος (Eur. Hec. 928) ; Synag.~κ~263 κελαδοῦσιν· ἠχοῦσιν ; \textit{idem} Suid.~κ~1281 (ἠχοῦσι ; add. σάλπιγξ ἄγκειμαι Τριτωνίδι κούρᾳ δῶρον, ἐριβρύχων παυσαμένη κελάδων. τουτέστι θορύβων) ; Suid.~κ~1279 Κελαδήματα: ἠχήματα ; \textit{Test.} Suid.~κ~1278 Κελαδείτω ; Ar.~\textit{Nub.}~311–313 ἦρί τ᾽ ἐπερχομένῳ Βρομία χάρις, / εὐκελάδων τε χορῶν ἐρεθίσματα, / καὶ μοῦσα βαρύβρομος αὐλῶν ; Ar.~\textit{Pax}~800–801 ὅταν ἠρινὰ μὲν φωνῇ χελιδὼν ἑζομένη κελαδῇ ; Ar.~\textit{Ran.}~384–385 ζαθέαις μολπαῖς κελαδεῖτε ; Ar.~\textit{Ran.}~683 κελαδεῖ δ᾽ ἐπίκλαυτον ἀηδόνιον νόμον ; Ar.~\textit{Ran.}~1525–1527 τοῖσιν τούτου τοῦτον μέλεσιν καὶ μολπαῖσιν κελαδοῦντες ; Eur.~\textit{HF}~678–679 ἔτι τοι γέρων ἀοιδὸς κελαδεῖ Μναμοσύναν ; Eur.~\textit{HF}~692–694 κύκνος ὣς γέρων ἀοι- / δὸς πολιᾶν ἐκ γενύων / κελαδήσω ; Aesch.~\textit{Cho.}~608–609 δαλὸν ἥλικ᾽, ἐπεὶ μολὼν / ματρόθεν κελάδησε ; \textit{Il.}~18.576 πὰρ ποταμὸν κελάδοντα ; Pratin. ap. Ath.~14.617c ἐμός, ἐμὸς ὁ Βρόμιος, ἐμὲ δεῖ κελαδεῖν, ἐμὲ δεῖ παταγεῖν ; Hsch.~κ~2125 κελαδεῖ· βοᾷ ; Hsch.~κ~2126 κελαδεινή· κραυγή, βοή. ἢ κυνηγὸς ἐπιθετικῶς (Υ 70) ; Hsch.~κ~2130 κελάδων· ἠχῶν, μέλπων (Φ 16). καὶ ποταμὸς Ἀρκαδίας (Η 133)}
\lp{45 γλαυκόν]}{Ar.~\textit{Av.}~1337–1339 γενοίμαν αἰετὸς ὑψιπέτας, ὡς ἀμποταθείην ὑπὲρ ἀτρυγέτου γλαυκᾶς ἐπ᾽ οἶδμα λίμνας ; Ar.~\textit{Ran.}~665 ὃς Αἰγαίου πρῶνας ἢ γλαυκᾶς μέδεις ἁλὸς ἐν βένθεσιν ; Eur.~\textit{Hel.}~400 ἐγὼ δ᾽ ἐπ᾽ οἶδμα πόντιον γλαυκῆς ἁλὸς ; Eur.~\textit{Hel.}~1501–1502 γλαυκὸν ἔπιτ᾽ οἶδμα κυανόχροά τε κυμάτων ῥόθια πολιὰ θαλάσσας ; \textit{Il.}~16.34 γλαυκὴ δέ σε τίκτε θάλασσα ; Σ Ar.~\textit{Ran.}~665 (= Soph. fr.~371~Radt) Παρὰ τὰ Σοφοκλέους ἐκ Λαοκόωντος « Πόσειδον, ὃς Αἰγαίου μέδεις πρῶνας, ἢ γλαυκᾶς μέδεις εὐανέμου λίμνας, ἐφ᾽ ὑψηλαῖς σπιλάδεσσι στομάτων. » ; Hsch.~γ~602 γλαυκή· ἰσχυρά. φοβερά. λευκή (Π 34) ; Eur.~\textit{Cyc.}~16 παῖδες δ᾽ ἐρετμοῖς ἥμενοι γλαυκὴν ἅλα ; Hes.~\textit{Th.}~440 καὶ τοῖς, οἳ γλαυκὴν δυσπέμφελον ἐργάζονται ; Hsch.~γ~606 γλαυκός· λευκός ; Synag.~γ~50 γλαυκός· λευκός, κυάνεος}
\lp{45 βομβάξ]}{\textit{Test.} Hsch.~β~787 βόμβαξ· τοῦτο παρεμβολοειδές ἐστι. σημαίνει δὲ διασυρμόν (Ar. Thesm. 45) ; Suid.~ε~2807 βομβάξ, ἐπίρρημα ἐπὶ θαυμασμοῦ λαμβανόμενον ; \textit{idem} Σ Ar.~\textit{Thesm.}~45 (om. βομβάξ) ; Eust.~1718.9–11 (ad Od.~12.204) δῆλον δέ, ὡς καὶ λίθος βομβεῖ … καὶ μέλισσα δὲ καὶ βροντή. ἐκ τούτου δὲ καὶ ἔντομόν τι ζῷον ὁ βομβυλιός, καὶ παρὰ τῷ κωμικῷ δὲ βομβὰξ παικτικὸν ἐπίρρημα καθ᾽ ὁμοιότητα τοῦ βαβαί. ὡσαύτως καὶ βομβαλοβομβάξ ; Ar.~\textit{Ach.}~64 σίγα. — βαβαιάξ. ὦκβάτανα τοῦ σχήματος ; Ar.~\textit{Ran.}~1285–1289 ὅπως Ἀχαιῶν δίθρονον κράτος, Ἑλλάδος ἥβας, τοφλαττοθρατ τοφλαττοθρατ, Σφίγγα δυσαμεριᾶν πρύτανιν κύνα, πέμπει, τοφλαττοθρατ τοφλαττοθρατ ; Plaut.~\textit{Pseud.}~365 Ballio : Babae. Pseudolus : Fugitive. Ballio : Bombax ; Pl.~\textit{Euthd.}~303a ὁ δὲ Κτήσιππός μοι ἰὼν ὡς βοηθήσων, Πυππὰξ ὦ Ἡράκλεις, ἔφη, καλοῦ λόγου ; Hsch.~κ~3184 κόγξ· ὁμ⟨οίως⟩ πάξ. ἐπιφώνημα τετελεσμένοις ; Ar.~\textit{Pax}~248 βαβαὶ βαβαιὰξ ὡς μεγάλα καὶ δριμέα ; Ar.~\textit{Thesm.}~945 ἰαππαπαιάξ· ὦ κροκώθ᾽ οἷ᾽ εἴργασαι ; Diph.~fr.~96~K.-A. ap. Ath.~2.67d ὄξους δὲ κοτύλην β. πάξ. α. τί πάξ ; Aesch.~\textit{Eum.}~143 ἰοὺ ἰοὺ πύπαξ. ἐπάθομεν, φίλαι ; Hsch.~β~2 βαβαί· θαυμαστικὴ φωνή ; Hsch.~β~9 βάβαξ· μάταιος, λάλος, φλύαρος (Archil. fr. 33). ἐνθουσιῶν. ἀναιδής ; Hsch.~π~433 πάξ· ὑπόδημα εὐυπόδητον (Herond. 7, 114). ἢ τέλος ἔχει[ν] (Diph. com. fr. 96, 2 K.–A. Men. Epitr. 657 Koerte = 987 Sandbach) ; Hsch.~π~4378 πυππάζουσι· φωνῇ ποιᾷ χρῶνται (Cratin. fr. 56 K.–A.) ; Hsch.~π~4379 πύππαξ· τὸ νῦν βόμβαξ λεγόμενον πύππαξ ἔλεγον, ὡς Λυκόφρων (fr. 115 Strecker) ᾠήθη. οὐκ ἔστι δέ. τὸ μὲν γὰρ βόμβαξ τίθεται καὶ ἐπὶ σχετλιασμοῦ καὶ ἐπὶ γέλωτος, τὸ δὲ πύππαξ οὐχί (Plat. Euthyd. 303 a. Aesch. Eum. 143. com. ad. 417 K.–A.?) ; Hsch.~φ~1019 φύππαξ· ὅπερ ἡμεῖς βόμβαξ ; Phot.~π~1543 πύππαξ· ἐπίφθεγμα σχετλιασμοῦ· ὡς πένθους ἀμετάφραστον· ὡς τὸ τρισκαιδέκατον ; Suid.~β~370 Βομβάζειν: φθέγμα ἐπὶ τῷ διασύρειν καὶ τωθάζειν λεγόμενον ; Ar.~\textit{Vesp.}~235 ἀππαπαῖ παπαιάξ ; Ar.~\textit{Av.}~1258 οὐκ ἀποσοβήσεις; οὐ ταχέως; εὐρὰξ πατάξ ; Pl. Com. ap. Ath.~15.666d βαβαιάξ· οὑτοσὶ μείζων ἀγὼν τῆς Ἰσθμιάδος ἐπέρχεται ; Eur.~\textit{Cyc.}~153 παπαιάξ, ὡς καλὴν ὀσμὴν ἔχει}
\lp{45 βομβάξ (βόμβος)]}{Ar.~\textit{Ach.}~864–866 οἱ σφῆκες οὐκ ἀπὸ τῶν θυρῶν; πόθεν προσέπτονθ᾽ οἱ κακῶς ἀπολούμενοι ἐπὶ τὴν θύραν μοι Χαιριδῆς βομβαύλιοι ; Ar.~\textit{Vesp.}~106–108 ὥσπερ μέλιττ᾽ ἢ βομβυλιὸς εἰσέρχεται ὑπὸ τοῖς ὄνυξι κηρὸν ἀναπεπλασμένος ; Ar.~\textit{Av.}~748–750 ἔνθεν ὡστερεὶ μέλιττα Φρύνιχος ἀμβροσίων μελέων ἀπεβόσκετο καρπὸν ἀεὶ φέρων γλυκεῖαν ᾠδάν ; Ar.~\textit{Thesm.}~1176 τί τὸ βόμβο τοῦτο; κῶμο τίς ἀνεγεῖρί μοι ; Aesch.~fr.~57~N.² (\textit{Edoni}) ap. Str.~10.3.16 ὁ μὲν ἐν χερσὶν βόμβυκας ἔχων, τόρνου κάματον, δακτυλόδικτον πίμπλησι μέλος ; Theoc.~5.29 σφὰξ βομβέων τέττιγος ἐναντίον ; \textit{idem} Apostol.~15.85b (σφῆξ βομβῶν ; add. ἐπὶ τῶν ἐριζόντων τοῖς κρείττοσι τάττεται) ; Pl.~\textit{Prt.}~316a διὰ τὴν βαρύτητα τῆς φωνῆς βόμβος τις ἐν τῷ οἰκήματι γιγνόμενος ἀσαφῆ ἐποίει τὰ λεγόμενα ; Pl.~\textit{R.}~8.564d τὸ δ᾽ ἄλλο περὶ τὰ βήματα προσίζον βομβεῖ τε καὶ οὐκ ἀνέχεται τοῦ ἄλλα λέγοντος ; Hsch.~σ~2886 σφηκισμός· εἶδος αὐλήσεως εἰρημένον ἀπὸ τῆς ἐμφερείας τῶν βόμβων ; Suid.~β~372 Βόμβος δὲ ὁ ἦχος τῶν μελισσῶν. βόμβος, ψόφος τις ; Suid.~β~375 Βομβυλιός: ζῷον, ἢ τὸ βησίον λεγόμενον. καὶ εἶδος μελίσσης, παρὰ τὸ βομβεῖν. ἐκ πηλοῦ τὰ κηρία πλαττούσης ; Zen.~2.80 Βομβύλιος ἄνθρωπος; ἐπὶ τοῦ ἀκάρπου· παρόσον ὁ βομβύλιος καρπὸν οὐ φέρει. Ἔστι δὲ μελίττης εἶδος ἐκ πηλοῦ τὰ κηρία πλαττούσης ; \textit{idem} Suid.~β~376 (Βομβυλιὸς ἄνθρωπος ; om. παρόσον … πλαττούσης) ; Zen.~3.53 Εἰς μελίττας ἐκώμασας: παροιμία ἐπὶ τῶν παρὰ δόξαν κακουμένων ἀθρόως ; \textit{idem} App. Prov.~2.27 (ἐκώμασεν ; om. παροιμία et παρὰ δόξαν) ; Ar.~\textit{Lys.}~475 ἢν μή τις ὥσπερ σφηκιὰν βλίττῃ με κἀρεθίζῃ ; Ar.~\textit{Plut.}~538–539 αἳ βομβοῦσαι περὶ τὴν κεφαλὴν ἀνιῶσιν, ἐπεγείρουσαι ; Pl.~\textit{Cri.}~54d καὶ ἐν ἐμοὶ αὕτη ἡ ἠχὴ τούτων τῶν λόγων βομβεῖ καὶ ποιεῖ μὴ δύνασθαι τῶν ἄλλων ἀκούειν ; Hsch.~β~788 βομβαύλιος· ὁ αὐλητής. ἀπὸ τοῦ βομβεῖν (Ar. Ach. 866) ; Hsch.~β~793 βόμβους· ἤχους, ψόφους ; Hsch.~β~800 βόμβυκες· γένος αὐλῶν, ἢ εἶδος ζώου πτερωτοῦ, κατὰ σφῆκα ; Hsch.~β~802 βομβυλιός· … ἢ ζῶον, ἦχόν τινα ποιοῦν τοῦ γένους τῶν σφηκῶν (Isocr. 10, 12). ἢ μέλισσα μεγάλη. ἢ μυῖα ; Hsch.~β~803 βόμβυξ· στάμνος. Λάκωνες. καὶ αὐλοῦ εἶδος. καὶ ζῶον … ἢ ἦχος βομβυλισμοῦ ; Suid.~β~374 Βομβοῦσιν αἱ μέλισσαι, ὅταν τις προσίῃ μύρου ὄζων ; Erasm.~\textit{Adag.}~2.6.71 Homo Bombylius. Βομβύλιος ἄνθρωπος. i. e. Bombylius homo, dicitur verbosus, multique strepitus, caeterum inutilis ; Plaut.~\textit{Amph.}~707 Inritabis crabrones ; Diog. Ath. ap. Ath.~14.636a τυπάνοισι καὶ ῥόμβοισι καὶ χαλκοκτύπων βόμβοις βρεμούσας ἀντίχερσι κυμβάλων ; \textit{Od.}~13.106 ἔνθα δ᾽ ἔπειτα τιθαιβώσσουσι μέλισσαι ; Pl.~\textit{R.}~9.573a ὅταν δὴ περὶ αὐτὸν βομβοῦσαι αἱ ἄλλαι ἐπιθυμίαι ; Plu.~\textit{Praec.\,ger.\,reip.}~32 τὸν μάλιστα βομβοῦντα τῶν σίμβλων καὶ θορύβου μεστὸν ; Hsch.~β~789 βομβῆσαι· ὡς περιστερὰ φωνῆσαι ; Hsch.~β~790 βόμβησεν· ἐψόφησεν (Ν 530 . .) ; Hsch.~β~792 βόμβος· ἦχος, ψόφος (Plat. Prot. 316a) ; Hsch.~β~796 βομβρύζων· τονθρύζων. βοῶν ; Hsch.~β~799 βομβύλην· λήκυθον ; Hsch.~β~1404 βωσσοῦσι· βομβοῦσιν, ἠχοῦσι (ν 106) ; Hsch.~υ~548 ὑπηχεῖ· βομβεῖ (Plat. Phaedr. 230 c) ; Harp. s.v.~βομβυλιός ΒΟΜΒΥΛΙΟΣ ζῷον παραπλήσιον μελίττῃ, ὠνομασμένον ἀπὸ τοῦ βόμβου· Ἰσοκράτης Ἑλένης ἐγκωμίῳ ; Diogenian.~1.15 Αὐλὸν σάλπιγγι συγκρίνεις: ἐπὶ τῶν τὰ ἐλάττω τοῖς μείζοσι συγκρινόντων. Τέττιγι μέλιτταν συγκρίνεις, κατὰ φωνήν}
\lp{45–48 βομβάξ … βομβαλοβομβάξ]}{Ar.~\textit{Nub.}~293–297 καὶ σέβομαί γ᾽ ὦ πολυτίμητοι καὶ βούλομαι ἀνταποπαρδεῖν πρὸς τὰς βροντάς … οὐ μὴ σκώψει μηδὲ ποιήσεις ἅπερ οἱ τρυγοδαίμονες οὗτοι, ἀλλ᾽ εὐφήμει· μέγα γάρ τι θεῶν κινεῖται σμῆνος ἀοιδαῖς}
\lp{45 σίγα. τί λέγει;]}{Ar.~\textit{Ach.}~238 σῖγα πᾶς. ἠκούσατ᾽ ἄνδρες ἆρα τῆς εὐφημίας ; Ar.~\textit{Lys.}~769 λέγ᾽ αὐτὸν ἡμῖν ὅ τι λέγει. — σιγᾶτε δή ; Ar.~\textit{Thesm.}~27 σίγα νυν. — σιωπῶ τὸ θύριον ; Ar.~\textit{Thesm.}~95 σίγα. — τί δ᾽ ἔστιν; — Ἁγάθων ἐξέρχεται ; Plaut.~\textit{Mil.}~993 Tace, subauscultemus ecquid de me fiat mentio ; Hsch.~ι~141 ἵγα· σιώπα. Κύπριοι ; Hsch.~σ~473 σιγᾷ· σιωπᾷ ; Plu.~\textit{De\,garr.}~4 (503d–504c) ὁ μὲν Βίας ἔν τινι πότῳ σιωπῶν καὶ σκωπτόμενος εἰς ἀβελτερίαν ὑπό τινος ἀδολέσχου, καὶ τίς ἄν ἔφη δύναιτο μωρὸς ὢν ἐν οἴνῳ σιωπᾶν ; Apostol.~11.92 Μωρὸς σιωπᾶν οὐ δύναται ; App. Prov.~3.7 Ἢ λέγε τι σιγῆς κρεῖττον, ἢ σιγὴν ἔχε: ἐπὶ τῶν οὐκ ἄξια λεγόντων ; Tosi n°~2072 Τὸ σιγᾶν πολλάκις ἐστὶ σοφώτατον ἀνθρώπῳ νοῆσαι. Souvent se taire est la plus sage des pensées}
\lp{46 πτηνῶν τε γένη]}{Ar.~\textit{Av.}~1063–1064 κτείνων παμφύλων γένναν θηρῶν ; Ar.~\textit{Av.}~1088 εὔδαιμον φῦλον πτηνῶν ; Ar.~\textit{Av.}~1706–1707 ὦ πάντ᾽ ἀγαθὰ πράττοντες, ὦ μείζω λόγου, ὦ τρισμακάριον πτηνὸν ὀρνίθων γένος ; Anaxil. ap. Ath.~13.558a Σφίγξ, Ὕδρα, λέαιν᾽, ἔχιδνα, πτηνὰ θ᾽ Ἁρπυιῶν γένη ; Hsch.~π~4219 πτηνός· ὁ ἱπτάμενος (Eur. Hipp. 1292. Ion 1196. Ar. Pac. 126) ; Ar.~\textit{Pax}~126 πτηνὸς πορεύσει πῶλος· οὐ ναυσθλώσομαι ; Ar.~\textit{Pax}~141 πῶς ἐξολισθεῖν πτηνὸς ὢν δυνήσεται ; Eur.~\textit{Andr.}~727 ἀνειμένον τι χρῆμα πρεσβυτῶν γένος ; Eur.~\textit{Hipp.}~1252 οὐδ᾽ εἰ γυναικῶν πᾶν κρεμασθείη γένος ; Aesch.~\textit{Eum.}~912 τὸ τῶν δικαίων τῶνδ᾽ ἀπένθητον γένος ; Soph.~\textit{Aj.}~397 οὔτε γὰρ θεῶν γένος οὔθ᾽ ἁμερίων}
\lp{46 κατακοιμάσθω]}{Pi.~\textit{P.}~1.6–9 εὕδει δ᾽ ἀνὰ σκάπτῳ Διὸς αἰετός, ὠκεῖαν πτέρυγ᾽ ἀμφοτέρωθεν χαλάξαις, ἀρχὸς οἰωνῶν ; Hsch.~κ~1719 κατευνάζεσθαι· κοιμίζεσθαι. πραΰνεσθαι (Cyr. in xii proph. 1, 163, 9) ; Hsch.~π~3225 πραΰνει· κατασιγαίνει. κατακοιμίζει ; \textit{Il.}~9.427 Φοῖνιξ δ᾽ αὖθι παρ᾽ ἄμμι μένων κατακοιμηθήτω ; Hsch.~β~360 βαυκαλᾶν· κατακοιμίζειν. τιθηνεῖν ⟨τὰ⟩ παιδία μετ᾽ ᾠδῆς, κοιμίζειν ; Hsch.~κ~1172 κατακοιμίζει· καταπραΰνει (Cyr. in Ioh. 1, 536, 12) ; Hsch.~κ~1718 κατευνάζει· κοιμίζει (Soph. Tr. 95) ; Hdt.~1.31.5 κατακοιμηθέντες ἐν αὐτῷ τῷ ἱρῷ οἱ νεηνίαι οὐκέτι ἀνέστησαν ; Hsch.~δ~322 δαύειν· κοιμᾶσθαι (Sapph. fr. 126 V.) ; Hsch.~κ~3250 κοιμηθέντι· κατακλιθέντι, οὐ πάντως ὑπνώσαντι (υ 4) ; Hsch.~κ~3251 κοιμηθῆναι· κατακλιθῆναι. ὑπνῶσαι (η 343) ; Hsch.~κ~3252 κοιμίζει· κατακλίνει ; Moer. p.~192~Bekker βαυκαλᾶν Ἀττικοί, κατακοιμίζειν Ἕλληνες ; Suid.~κ~1013 Κατευνάζεται: κατακοιμίζεται ; Zen.~3.76 Ἐνδυμίωνος ὕπνος: ἐπὶ τῶν πολλὰ κοιμωμένων ; Macar.~6.69 Οὐδ᾽ ἀηδόνιον κατέδαρθον: πρὸς τοὺς φιλαγρύπνους}
\lp{46–47 πτηνῶν … θηρῶν]}{Soph.~\textit{Ant.}~343–345 κουφονόων τε φῦλον ὀρνίθων ἀμφιβαλὼν ἄγει καὶ θηρῶν ἀγρίων ἔθνη πόντου τ᾽ εἰναλίαν φύσιν ; Verg.~\textit{Aen.}~8.26–27 Nox erat, et terras animalia fessa per omnis / alituum pecudumque genus sopor altus habebat ; Eur.~\textit{Ba.}~561–564 ἔνθα ποτ᾽ Ὀρφεὺς κιθαρίζων σύναγεν δένδρεα μούσαις, σύναγεν θῆρας ἀγρώτας ; Hes.~\textit{Op.}~277 ἰχθύσι μὲν καὶ θηρσὶ καὶ οἰωνοῖς πετεηνοῖς ; \textit{h.Ven.}~4–5 οἰωνούς τε διιπετέας καὶ θηρία πάντα, ἠμὲν ὅσ᾽ ἤπειρος πολλὰ τρέφει ἠδ᾽ ὅσα πόντος ; Palaeph.~33 ὅτι κιθαρίζοντι αὐτῷ ἐφείπετο τετράποδα καὶ ἑρπετὰ καὶ ὄρνεα καὶ δένδρα}
\lp{47 θηρῶν τ᾽ ἀγρίων]}{Eur.~\textit{Alc.}~578–586 σὺν δ᾽ ἐποιμαίνοντο χαρᾷ μελέων βαλιαί τε λύγκες ; Soph.~\textit{Ant.}~348–350 κρατεῖ δὲ μηχαναῖς ἀγραύλου θηρὸς ὀρεσσιβάτα ; Hsch.~α~801 ἄγριον· ἄμικτον. ἀνήμερον (μ 119) ; Hsch.~α~805 ἄγριον· γενικῶς ἐπὶ θηρίου· «σῦν ἄγριον ἀργιόδοντα» (Ι 539) ; Hsch.~θ~515 θήρ· θηρίον. ὁ λέων· ἀπὸ τοῦ θοῶς ὀρούειν (Γ 449 . .) ; Hsch.~θ~524 θηρία· τὰ ζῷα. καὶ ἐπὶ τῶν ἵππων οὕτω λέγουσι}
\lp{47 ὑλοδρόμων]}{Ar.~\textit{Av.}~1393–1395 εἴδωλα πετεινῶν αἰθεροδρόμων οἰωνῶν ταναοδείρων ; Ael.~\textit{NA}~6.26 ἡ πιθήκη ὑπό τινων ὀρειβάτης κέκληται, ὑπ᾽ ἄλλων γε μήν, ὡς ἀκούω, ὑλοδρόμος ; Hsch.~υ~171 ὑλομανής· ὁ ταῖς ὕλαις χαίρων (Cyr. in Ep. Hebr. p. 406, 13 . .) ; Hsch.~ο~1162 ὀρέστερος· ἄγριος (Χ 93 . .) ; Eur.~\textit{Andr.}~284 ταὶ δ᾽ ἐπεὶ ὑλόκομον νάπος ἤλυθον οὐρειᾶν ; Eur.~\textit{IA}~207 λαιψηροδρόμον Ἀχιλῆα ; Eur.~\textit{IT}~261 βοῦς ὑλοφορβοὺς πόντον εἰσεβάλλομεν ; Catull.~63.72 ubi cerva silvicultrix, ubi aper nemorivagus? ; Hsch.~υ~154 ὑλῆεν· δασύ, ξυλῶδες (Nic. ther. 147. Simon. AP 6, 217, 8 = HE 3311) ; Hsch.~υ~155 ὑληέσσης· συμφύτου, ἢ ὕλας ἐχούσης (Ν 12 . .) ; Hsch.~υ~157 ὑληρεύς· νομεὺς ἐν ὕλῃ φυλάττων}
\lp{47–48 πόδες … μὴ λυέσθων]}{Ar.~\textit{Ran.}~1323 ὁρᾷς τὸν πόδα τοῦτον; — ὁρῶ ; Eur.~\textit{Med.}~679 ἀσκοῦ με τὸν προύχοντα μὴ λῦσαι πόδα— ; Soph.~\textit{Aj.}~675–676 ἐν δ᾽ ὁ παγκρατὴς ὕπνος λύει πεδήσας, οὐδ᾽ ἀεὶ λαβὼν ἔχει ; Σ Eur.~\textit{Med.}~679 ἀσκοῦ οὖν τῆς γαστρός· πόδα δὲ τὸ μόριον, παρόσον ὡς ὁ ποδεὼν τοῦ ἀσκοῦ προέχει. … ποδεῶνα δὲ εἰώθασι λέγειν τὸ τοῦ ἀνδρὸς αἰδοῖον ; Epicr. ap. Ath.~11.782f κατάβαλλε τἀκάτια, τὰ κυμβία αἴρου τὰ μείζω … εὐτρεπῆ τε τὸν κοντὸν ποοῦ καὶ τοὺς κάλως ἔκλυε καὶ χάλα πόδα ; Eur.~\textit{Ba.}~447 αὐτόματα δ᾽ αὐταῖς δεσμὰ διελύθη ποδῶν ; Eur.~\textit{Hec.}~1019–1020 καὶ γὰρ Ἀργεῖοι νεῶν λῦσαι ποθοῦσιν οἴκαδ᾽ ἐκ Τροίας πόδα ; Plu.~\textit{Thes.}~3.3 ἀσκοῦ τὸν προὔχοντα πόδα, μέγα φέρτατε λαῶν, μὴ λύσῃς πρὶν δῆμον Ἀθηνέων εἰσαφικέσθαι ; Phot.~π~982 ποδεών· κυρίως τοῦ ἀσκοῦ τὰ προὔχοντα· ἤτοι τῶν ποδῶν τὰ δέρματα ; Eur.~\textit{Or.}~706–707 καὶ ναῦς γὰρ ἐνταθεῖσα πρὸς βίαν ποδὶ ἔβαψεν, ἔστη δ᾽ αὖθις, ἢν χαλᾷ πόδα}
\lp{48 βομβαλοβομβάξ]}{\textit{Test.} Suid.~ε~2807 βομβαλομβάξ, ἀναδιπλασιάζει τὴν λέξιν ; \textit{idem} Σ Ar.~\textit{Thesm.}~48 (om. βομβαλομβάξ) ; Eust.~790.13 (ad Il.~10.67) δι᾽ ὃν ποιεῖ βομβώδη κόμπον, οὗ περ ἀκούσας ὁ κωμικὸς ἔπαιξεν ἄν, οἶμαι, καὶ εἰς αὐτὸν τὸ βομβαλοβομβάξ ; Ar.~\textit{Nub.}~390–391 ἀτρέμας πρῶτον παππὰξ παππάξ, κἄπειτ᾽ ἐπάγει παπαπαππάξ, χὤταν χέζω, κομιδῇ βροντᾷ παπαπαππὰξ ὥσπερ ἐκεῖναι ; \textit{Test.} Eust.~1718.10–11 ὡσαύτως καὶ βομβαλοβομβάξ ; Ar.~\textit{Ran.}~209–210 βρεκεκεκὲξ κοὰξ κοάξ, βρεκεκεκὲξ κοὰξ κοάξ ; Hsch.~β~177 βαμβαλεῖν· τρέμειν. ψοφεῖν τοῖς χείλεσι ; Suid.~π~257 Παππάξ παπαππάξ: ἦχος πορδῆς. Ἀριστοφάνης Νεφέλαις ; Ar.~\textit{Av.}~260–262 τοροτοροτοροτοροτίξ. κικκαβαῦ κικκαβαῦ. τοροτοροτοροτορολιλιλίξ ; Ar.~\textit{Thesm.}~223 ἢν μὴ σιωπᾷς. — ἀτταταῖ ἰατταταῖ ; Hsch.~β~180 βαμβαλύζει· τρέμει, τοὺς ὀδόντας συγκρούει. ῥιγοῖ σφόδρα}
\lp{48 μὴ λυέσθων]}{\textit{Il.}~7.6 καμάτῳ δ᾽ ὑπὸ γυῖα λέλυνται ; \textit{Od.}~20.56–57 εὖτε τὸν ὕπνος ἔμαρπτε, λύων μελεδήματα θυμοῦ, λυσιμελής ; Hsch.~λ~1445 λυσιμελής· ἐπιθετικῶς, ὁ ὕπνος. ἤτοι ὁ τὰ μέλη τοῦ σώματος λύων, ἢ τὰς μελεδώνας, τὰς μερίμνας λύων, «μελεδήματα θυμοῦ» (υ 57) ; Ar.~\textit{Nub.}~1142 νῦν οὖν δικαζέσθων· ὀλίγον γάρ μοι μέλει ; Hsch.~λ~634 λέλυνται· ἀσθενῆ εἰσιν (Β 135)}
\lp{49 ὁ καλλιεπὴς Ἀγάθων]}{\textit{Test.} Philostr.~\textit{VS}~1.9.1 καὶ Ἀγάθων δὲ ὁ τῆς τραγῳδίας ποιητής, ὃν ἡ κωμῳδία σοφόν τε καὶ καλλιεπῆ οἶδε, πολλαχοῦ τῶν ἰάμβων γοργιάζει ; Hsch.~α~281 Ἀγαθώνειον αὔλησιν τὴν μαλακὴν (Ar. fr. 178 K.–A.). Ἀγαθὼν γὰρ ὁ τραγικὸς ἐπὶ μαλακίᾳ διεβάλλετο (TrGF 39 T 20a) ; Suid.~α~124 Ἀγάθων· ὄνομα κύριον. τραγικὸς δὲ ἦν· διεβέβλητο δὲ ἐπὶ μαλακίᾳ ; Zen.~1.2 Ἀγαθώνιος αὔλησις: ἡ μαλακὴ, καὶ μήτε πικρὰ μήτε χαλαρὰ, ἀλλ᾽ εὔκρατος καὶ ἡδίστη. Εἴρηται δὲ ἀπὸ Ἀγάθωνος αὐλητοῦ, ὃς ἐπὶ μαλακίᾳ ἐκωμῳδεῖτο ; \textit{idem} Diogenian.~1.7 (Ἀγαθώνειος ; μαλθακὴ καὶ μήτε χλιαρὰ μήτε πικρὰ ; om. εὔκρατος καὶ et Εἴρηται … ἐκωμῳδεῖτο), Suid.~α~125 (ἡ μαλακὴ καὶ ἐκλελυμένη· ἢ ἡ μήτε χαλαρὰ, μήτε πικρὰ ; om. Εἴρηται … ἐκωμῳδεῖτο) ; Erasm.~\textit{Adag.}~2.4.63 Agathonia cantio. Ἀγαθώνειος αὔλησις. i. e. Agathonia cantio. Recte dicetur de oratione blanda magis, quam frugifera. … Laboravit enim infamia molliciei. … Hujus multis modis Aristophanes in Thesmophoriis carminum molliciem deridet ; Pl.~\textit{Smp.}~174b ἵνα καὶ τὴν παροιμίαν διαφθείρωμεν μεταβαλόντες, ὡς ἄρα καὶ Ἀγάθων᾽ ἐπὶ δαῖτας ἴασιν αὐτόματοι ἀγαθοί ; Ar.~\textit{Ran.}~83–84 Ἀγάθων δὲ ποῦ ᾽στιν; — ἀπολιπών μ᾽ ἀποίχεται, ἀγαθὸς ποιητὴς καὶ ποθεινὸς τοῖς φίλοις ; Arist.~\textit{Po.}~18.20 διὸ ἐμβόλιμα ᾄδουσιν πρώτου ἄρξαντος Ἀγάθωνος τοῦ τοιούτου ; Zen.~2.19 Αὐτόματοι δ᾽ ἀγαθοὶ ἀγαθῶν ἐπὶ δαῖτας ἴενται ; App. Prov.~3.43 Καπνίας ποιητής: ὁ μηδὲν λαμπρὸν γράφων ; Suid.~α~123 Ἀγαθῶν ἀγαθίδες: τάττεται ἡ παροιμία παρὰ τοῖς κωμικοῖς ἐπὶ τῶν πολλῶν ἀγαθῶν}
\lp{49 καλλιεπής]}{Pl.~\textit{Smp.}~198b τὸ δὲ ἐπὶ τελευτῆς τοῦ κάλλους τῶν ὀνομάτων καὶ ῥημάτων τίς οὐκ ἂν ἐξεπλάγη ἀκούων ; Arist.~\textit{Rh.}~3.2.3 ἐπεὶ καὶ ἐνταῦθα, εἰ δοῦλος καλλιεποῖτο ἢ λίαν νέος, ἀπρεπέστερον ; Th.~6.83.2 καὶ οὐ καλλιεπούμεθα ὡς ἢ τὸν βάρβαρον μόνοι καθελόντες εἰκότως ἄρχομεν ; Pl.~\textit{Ap.}~17b–c οὐ μέντοι μὰ Δία, ὦ ἄνδρες Ἀθηναῖοι, κεκαλλιεπημένους γε λόγους, ὥσπερ οἱ τούτων, ῥήμασί τε καὶ ὀνόμασιν οὐδὲ κεκοσμημένους ; Pl.~\textit{Smp.}~201c κινδυνεύω, ὦ Σώκρατες, οὐδὲν εἰδέναι ὧν τότε εἶπον. καὶ μὴν καλῶς γε εἶπες, φάναι, ὦ Ἀγάθων ; Hsch.~κ~473 καλλιέπεια· καλλιλεξία ; Pl.~\textit{Hipparch.}~225c ἵνα τι καὶ ἡμεῖς τῶν σοφῶν ῥημάτων ἐμβάλωμεν, ὧν οἱ δεξιοὶ περὶ τὰς δίκας καλλιεποῦνται ; D.H.~\textit{Pomp.}~2.5 ὅταν δ᾽ εἰς τὴν περιττολογίαν καὶ τὸ καλλιεπεῖν ; Hsch.~ε~6792 εὐεπής· εὔλαλος (Cyr. Glaph. in Pent. 69, 481A) ; Hsch.~ε~7108 εὔρημον· καλλιεπές}
\lp{49 καλλιεπής (ἡδυεπής, καλλι-)]}{Pi.~\textit{N.}~7.20–21 λόγον Ὀδυσσέος ἢ πάθαν διὰ τὸν ἁδυεπῆ γενέσθ᾽ Ὅμηρον ; Ar.~\textit{Ach.}~1227 τήνελλα καλλίνικος ; \textit{idem} Ar.~\textit{Av.}~1764 ; Ar.~\textit{Eq.}~1254 ὦ χαῖρε καλλίνικε ; Ar.~\textit{Pax}~797 τοιάδε χρὴ Χαρίτων δαμώματα καλλικόμων ; Ar.~\textit{Av.}~682 ἀλλ᾽ ὦ καλλιβόαν κρέκουσ᾽ ; Ar.~\textit{Ran.}~452 τὸν καλλιχορώτατον ; Eur.~\textit{HF}~368–369 ξύνοιδε Πηνειὸς ὁ καλ- / λιδίνας ; Eur.~\textit{Hec.}~467 τὰς καλλιδίφρους ; Eur.~\textit{Ion}~189 καλλιβλέφαρον φῶς ; Eur.~\textit{Or.}~1382 Φρύγιον ἄστυ καὶ καλλίβωλον ; Eur.~\textit{Phoen.}~645 καλλιπόταμος ὕδατος ; Eur.~\textit{Phoen.}~1162 τῇ καλλιτόξῳ μητρὶ Μαινάλου κόρῃ ; Soph.~\textit{OT}~151 ὦ Διὸς ἁδυεπὲς φάτι ; \textit{Il.}~1.248 ἡδυεπὴς ἀνόρουσε λιγὺς Πυλίων ἀγορητής ; Pi.~\textit{O.}~10.93 τὶν δ᾽ ἁδυεπής τε λύρα}
\lp{50 πρανοῦν]}{\textit{Test.} Hsch.~π~3188 πραῖνοι· πρηνίζειν. καταστρέφειν ; Ar.~\textit{Ach.}~399–400 αὐτὸς δ᾽ ἔνδον ἀναβάδην ποιεῖ τραγῳδίαν ; Ar.~\textit{Pax}~896–897 ἐπὶ γῆς παλαίειν, τετραποδηδὸν ἑστάναι, πλαγίαν καταβάλλειν, ἐς γόνατα κύβδ᾽ ἑστάναι ; Ar.~\textit{Thesm.}~489 παρὰ τὸν Ἀγυιᾶ κύβδ᾽ ἐχομένη τῆς δάφνης ; Hsch.~ε~5537 ἐπράνιξε· κατέβαλεν ; Hsch.~π~3267 πρηνής· ἐπὶ πρόσωπον πεπτωκώς (Ε 58 . .)· ἀπὸ τοῦ ἐπὶ τὴν νειὸν πεπτωκέναι, νειὸς δὲ ἡ γῆ ; Ar.~\textit{Eq.}~365 ἐγὼ δέ γ᾽ ἐξέλξω σε τῆς πυγῆς θύραζε κύβδα ; Ar.~\textit{Thesm.}~97–98 ἀλλ᾽ ἢ τυφλὸς μέν εἰμ᾽· ἐγὼ γὰρ οὐχ ὁρῶ ἄνδρ᾽ οὐδέν᾽ ἐνθάδ᾽ ὄντα, Κυρήνην δ᾽ ὁρῶ ; Archil.~fr.~32~Bergk ap. Ath.~10.447a–d ὥσπερ παρ᾽ αὐλῷ βρῦτον ἢ Θρέιξ ἀνὴρ ἢ Φρὺξ ἔβρυζε· κύβδα δ᾽ ἦν πονευμένη ; Hsch.~κ~4656 κυπτάζειν· διαπρίβειν, στραγγεύειν. καὶ †ἐναισχύνειν. ἐπικύπτειν. καὶ χρονίζειν ; Hsch.~κ~4657 κύπτει· κάμπτει, ἑαυτὸν κλίνει ; Hsch.~π~3197 πρανές· κατάβασις (Xen. Cyr. 2, 2, 24. Anab. 3, 4, 26). καὶ πρανής (Xen. Anab. 5, 2, 29) ; Hsch.~π~3199 πρανιχθέντα· πεσόντα ἐπὶ στόμα ; Hsch.~π~3200 πρανόν· τὸ κατωφερές, πρανές ; Hsch.~π~3268 πρήνιξε· κατέβαλεν (Euphor. SH fr. 418, 41. Greg. Naz. c. 2, 1, 13, 54) ; Hsch.~π~3269 πρηνιχθέντα· ἐπὶ πρόσωπον καὶ ἐπὶ στόμα πεσόντα ; Hsch.~π~3603 προνωπής· προτεταμένος. προνενευκώς. οἱ δὲ προπετής, ἕτοιμος, πρόχειρος (Eur. Andr. 729) ; Phot.~π~1138 πρανιχθῆναι· τὸ ἐπὶ στόμα πεσεῖν}
\lp{50 προτέρους]}{Ar.~\textit{Nub.}~562 ἐς τὰς ὥρας τὰς ἑτέρας εὖ φρονεῖν δοκήσετε}
\lp{50 πρόμος (Grynaeus)]}{Aesch.~\textit{Supp.}~907 ἰὼ πόλεως ἀγοὶ πρόμοι, δάμναμαι ; Ar.~\textit{Nub.}~394 ταῦτ᾽ ἄρα καὶ τὠνόματ᾽ ἀλλήλοιν βροντὴ καὶ πορδὴ ὁμοίω ; Eur.~\textit{Heracl.}~670 Ἀθηναίων πρόμοι ; Aesch.~\textit{Ag.}~199–200 χείματος ἄλλο μῆχαρ / βριθύτερον πρόμοισιν ; Aesch.~\textit{Ag.}~410 ἰὼ ἰὼ δῶμα δῶμα καὶ πρόμοι ; Soph.~\textit{OT}~660 οὐ τὸν πάντων θεῶν θεὸν πρόμον ; \textit{Il.}~3.44 φάντες ἀριστῆα πρόμον ἔμμεναι, οὕνεκα καλὸν ; \textit{Od.}~11.493 ἢ ἕπετ᾽ ἐς πόλεμον πρόμος ἔμμεναι, ἦε καὶ οὐκί ; Hsch.~π~3585 πρόμοις· προμάχοις, προαγωνισταῖς, προστάταις, προηγουμένοις (Heraclid. Lemb. FHG 3, 169 fr. 5) ; Suid.~π~2516 Πρόμος: πρόμαχος}
\lp{50 μῶν βινεῖσθαι;]}{Ar.~\textit{Thesm.}~206 ἰδού γε κλέπτειν· νὴ Δία βινεῖσθαι μὲν οὖν ; Ar.~\textit{Ran.}~1208 Ἄργος κατασχών— — ληκύθιον ἀπώλεσεν ; Com. adesp. fr.~12–14~Kock οὐδεὶς κομήτης ὅστις οὐ ψηνίζεται. — οὐδεὶς κομήτης ὅστις οὐ βινητιᾷ. — οὐδεὶς κομήτης ὅστις οὐ περαίνεται ; \textit{idem} Phot.~ο~609 (οὐ περαίνεται), Suid.~κ~1996 (οὐ ψηνίζεται ; praem. οὐδεὶς κομήτης ὅστις οὔ. τὸ δὲ ἀκροτελεύτιον … φθέγξομαι τὸ δεινὸν ἐκεῖνο πρᾶγμα καὶ ὄνομα), Macar.~6.74 (οὐ βινητιᾷ ; add. ὡς ἐν διαβολῇ) ; Hsch.~β~466 βινεῖν· παρὰ Σόλωνι (legg. F 27) τὸ βίᾳ μίγνυσθαι. τὸ δὲ κατὰ νόμον ὀπύειν ; Ar.~\textit{Eq.}~877–879 ἔπαυσα τοὺς βινουμένους, τὸν Γρύττον ἐξαλείψας ; Ar.~\textit{Nub.}~1089–1094 συνηγοροῦσιν ἐκ τίνων; — ἐξ εὐρυπρώκτων. … τί δαί; τραγῳδοῦσ᾽ ἐκ τίνων; — ἐξ εὐρυπρώκτων. … δημηγοροῦσι δ᾽ ἐκ τίνων; — ἐξ εὐρυπρώκτων ; Ar.~\textit{Thesm.}~35 καὶ μὴν βεβίνηκας σύ γ᾽, ἀλλ᾽ οὐκ οἶσθ᾽ ἴσως ; Ar.~\textit{Eccl.}~112–113 λέγουσι γὰρ καὶ τῶν νεανίσκων ὅσοι — πλεῖστα σποδοῦνται, δεινοτάτους εἶναι λέγειν ; Hsch.~β~467 βινητιᾶν· τὸ πασχητιᾶν. καὶ τὸ ὄρεξιν ἐπὶ συνουσίαν ἔχειν (Plat. com. fr. 188, 21 K.–A.) ; Hsch.~σ~1534 σποδεῖν· [σ]παίειν (Ar. Pac. 1306. Av. 1016), συγγίνεσθαι (Ar. Eccl. 1016) ; Hsch.~σ~1535 σποδεῖσθαι· περαίνεσθαι (Ar. Eccl. 908) ; Hsch.~κ~3447 κομήτης· εὔθριξ, καλὰς τρίχας ἔχων ; Hsch.~μ~2057 μῶν· μὴ ⟨οὖν⟩. οὐ δή. ἆρα ; Hsch.~ο~432 οἰφεῖν· ὀχεύειν. καὶ τὰ ὅμοια ; Suid.~β~288 Βινεῖν: περαίνειν}
\lp{51 τίς ὁ φωνήσας;]}{Ar.~\textit{Av.}~60 ἐποποῖ. — τίνες οὗτοι; τίς ὁ βοῶν τὸν δεσπότην ; Ter.~\textit{An.}~267 Quis hic loquitur? Mysis? salve ; Eur.~\textit{IA}~857 τίς ὁ καλῶν πύλας παροίξας; ὡς τεταρβηκὼς καλεῖ ; Soph.~\textit{OT}~1300 τίς ὁ πηδήσας ; Plaut.~\textit{Amph.}~331 Certe enim hic nescio quis loquitur ; \textit{Il.}~1.201 καί μιν φωνήσας ἔπεα πτερόεντα προσηύδα ; \textit{Il.}~10.465 ὣς ἄρ᾽ ἐφώνησεν ; \textit{Od.}~4.370 ἡ δέ μευ ἄγχι στᾶσα ἔπος φάτο φώνησέν τε ; Hsch.~φ~1102 φωνήσας· εἰπών, λέξας (Α 201 . .)}
\lp{51 νήνεμος αἰθήρ]}{Ar.~\textit{Vesp.}~184–185 τίς εἶ ποτ᾽ ὦνθρωπ᾽ ἐτεόν; — Οὖτις νὴ Δία. — Οὖτις σύ; ποδαπός; — Ἴθακος Ἀποδρασιππίδου ; Eur.~\textit{Cyc.}~672–675 Οὖτίς μ᾽ ἀπώλεσ᾽. — οὐκ ἄρ᾽ οὐδεὶς ἠδίκει ; \textit{Od.}~9.366 Οὖτις ἐμοί γ᾽ ὄνομα· Οὖτιν δέ με κικλήσκουσι ; \textit{Od.}~9.408 ὦ φίλοι, Οὖτίς με κτείνει δόλῳ οὐδὲ βίηφιν ; Lib.~\textit{Decl.}~26.34 τὸ Δωδωναῖον ὑπερβαίνει χαλκεῖον. ἐκεῖνο μὲν γὰρ προσβαλλούσης αὐτῷ τῆς μάστιγος ὑπὸ τῶν ἀνέμων ἠχεῖ, νηνεμίας δὲ οὔσης σιγᾷ ; Zen.~6.5 Τὸ Δωδωναῖον χαλκεῖον: κεῖται παρὰ Μενάνδρῳ ἐν τῇ Ἀῤῥηφόρῳ. Εἴρηται δὲ ἐπὶ τῶν πολλὰ λαλούντων καὶ μὴ διαλειπόντων}
\lp{52–57 δρυόχους τιθέναι … καὶ λαικάζει]}{\textit{Test.} Suid.~δ~1547 Δρύοχοι: πάτταλοι, οἱ ἐντιθέμενοι ναυπηγουμένης νεώς. καταχρηστικῶς δὲ τὰ προοίμια. Ἀριστοφάνης· Ἀγάθων ἄρχεται δρυόχους τιθέναι δράματος ἀρχάς. κάμπτει δὲ νέας ἁψίδας ἐπῶν, τὰ δὲ τορεύει, τὰ δὲ κολλομελεῖ καὶ γνωμοτυπεῖ κἀντονομάζει καὶ κηροχυτεῖ καὶ γογγυλίζει καὶ χιανεύει καὶ λαικάζει ; \textit{idem} Suid.~χ~297 (Χιανεύει καὶ λαικάζει· Ἀριστοφάνης pro Δρύοχοι … Ἀριστοφάνης), Zonar. col.~569~Tittmann (νηός ; usque ad ἀρχάς ; add. ἢ δρυόχους τὰς πελέκεις κτλ.), Σ Ar.~\textit{Ach.}~79 (καὶ ἀλλαχοῦ· κάμπτει δὲ νέας ἀψῖδας ἐπῶν, γογγυλίζει καὶ λαικάζει), Suid.~λ~181 (ὁ αὐτός· κάμπτει δὲ νέας ἀψῖδας ἐπῶν, τὸ δὲ τορείη καὶ γογγυλίζει καὶ λαικάζει) ; Ar.~\textit{Ach.}~398–400 ὁ νοῦς μὲν ἔξω ξυλλέγων ἐπύλλια οὐκ ἔνδον, αὐτὸς δ᾽ ἔνδον ἀναβάδην ποιεῖ τραγῳδίαν ; Ar.~\textit{Eq.}~461–470 οὐκ ἐλάνθανεν τεκταινόμενα τὰ πράγματ᾽, ἀλλ᾽ ἠπιστάμην γομφούμεν᾽ αὐτὰ πάντα καὶ κολλώμενα. — καὶ ταῦτ᾽ ἐφ᾽ οἷσίν ἐστι συμφυσώμενα ἐγᾦδ᾽· ἐπὶ γὰρ τοῖς δεδεμένοις χαλκεύεται ; Ar.~\textit{Eq.}~529–530 ᾆσαι δ᾽ οὐκ ἦν ἐν ξυμποσίῳ πλὴν Δωροῖ συκοπέδιλε, καὶ τέκτονες εὐπαλάμων ὕμνων· οὕτως ἤνθησεν ἐκεῖνος ; Ar.~\textit{Ran.}~799–801 καὶ κανόνας ἐξοίσουσι καὶ πήχεις ἐπῶν καὶ πλαίσια ξύμπτυκτα— — πλινθεύσουσι ; Telecl.~fr.~39–40~Kock ap. Vit. Eur. Μνησίλοχός ἐστ᾽ ἐκεῖνος, ὃς φρύγει τι δρᾶμα καινὸν Εὑριπίδῃ, καὶ Σωκράτης τὰ φρύγαν᾽ ὑποτίθησιν ; Pi.~\textit{P.}~3.113 ἐξ ἐπέων κελαδεννῶν, τέκτονες οἷα σοφοὶ ; Pl.~\textit{Ion}~534b κοῦφον γὰρ χρῆμα ποιητής ἐστιν καὶ πτηνὸν καὶ ἱερόν, καὶ οὐ πρότερον οἷός τε ποιεῖν πρὶν ἂν ἔνθεός τε γένηται ; Zen.~3.7 Δαιδάλεια ποιήματα: οἱ μὲν ἐπὶ τῶν παραδόξων ἔργων, οἱ δὲ ἐπὶ τῶν ἀκριβούντων τὰς τέχνας φασὶν εἰρῆσθαι τὴν παροιμίαν ; App. Prov.~1.38 Ἀπὸ ἵππων ἐπ᾽ ὄνους: ἐπὶ τῶν ἀπὸ σεμνῶν εἰς εὐτελῆ μεταβαινόντων οἷον ἀπὸ γραμματικῆς εἰς χαλκευτικήν ; Ar.~\textit{Pax}~749–750 ἐποίησε τέχνην μεγάλην ἡμῖν κἀπύργωσ᾽ οἰκοδομήσας ἔπεσιν μεγάλοις καὶ διανοίαις ; Ar.~\textit{Ran.}~824 ῥήματα γομφοπαγῆ πινακηδὸν ἀποσπῶν ; Ar.~\textit{Ran.}~1004 ἀλλ᾽ ὦ πρῶτος τῶν Ἑλλήνων πυργώσας ῥήματα σεμνὰ ; Ar.~\textit{Plut.}~513 τίς χαλκεύειν ἢ ναυπηγεῖν ἢ ῥάπτειν ἢ τροχοποιεῖν ; Soph.~fr.~162~N.² ap. Poll.~7.117 βιαία γὰρ ἡ ἐν τῷ Σοφοκλέους Δαιδάλῳ «τεκτόναρχος μοῦσα» ; Pi.~\textit{N.}~3.4–5 μελιγαρύων τέκτονες κώμων νεανίαι ; Pl.~\textit{Phlb.}~56b–c τεκτονικὴν δέ γε … πῇ; κατά τε ναυπηγίαν καὶ κατ᾽ οἰκοδομίαν καὶ ἐν πολλοῖς ἄλλοις τῆς ξυλουργικῆς. κανόνι γὰρ οἶμαι καὶ τόρνῳ χρῆται καὶ διαβήτῃ καὶ στάθμῃ καί τινι προσαγωγίῳ κεκομψευμένῳ. … θῶμεν τοίνυν διχῇ τὰς λεγομένας τέχνας, τὰς μὲν μουσικῇ συνεπομένας ἐν τοῖς ἔργοις ἐλάττονος ἀκριβείας μετισχούσας, τὰς δὲ τεκτονικῇ πλείονος ; Democr. ap. Dio Chrys.~53.1 Ὅμηρος φύσεως λαχὼν θεαζούσης ἐπέων κόσμον ἐτεκτήνατο παντοίων ; Pl.~\textit{Ti.}~69a οἷα τέκτοσιν ἡμῖν ὕλη παράκειται ; Hsch.~τ~388 τεκτονουργός· ἀρχιτέκτων}
\lp{52 δρυόχους]}{\textit{Od.}~19.574 ἵστασχ᾽ ἑξείης, δρυόχους ὥς, δώδεκα πάντας ; A.R.~1.723–724 Παλλάς, ὅτε πρῶτον δρυόχους ἐπεβάλλετο νηὸς Ἀργοῦς, καὶ κανόνεσσι δάε ζυγὰ μετρήσασθαι ; Pl.~\textit{Ti.}~81b νέα μὲν οὖν σύστασις τοῦ παντὸς ζῴου, καινὰ τὰ τρίγωνα οἷον ἐκ δρυόχων ἔτι ἔχουσα τῶν γενῶν ; Σ Ar.~\textit{Thesm.}~52 προοίμια. κυρίως δὲ δρύοχοί εἰσιν οἱ ἐντιθέμενοι πάτταλοι ναυπηγουμένης νεώς ; App. Prov.~1.15 Ἅλις δὲ δρυόχων: ἐπὶ τῶν πολλάκις ἀρχομένων λόγων καὶ μὴ περαινόντων ; Epigr. ap. Ath.~5.209c πῶς δὲ κατὰ δρυόχων ἐπάγη σανίς, ἢ τίνι γόμφοι τμηθέντες πελέκει τοῦτ᾽ ἔκαμον τὸ κύτος ; Plu.~\textit{De\,fort.\,Rom.}~9 ὥσπερ ἐκ δρυόχων κραταιῶν συνιστάμενος ; Plb.~1.38.5 αὖθις ἔγνωσαν ἐκ δρυόχων εἴκοσι καὶ διακόσια ναυπηγεῖσθαι σκάφη ; Procop.~\textit{Goth.}~8.22.12 ἅπερ οἱ μὲν ποιηταὶ δρυόχους καλοῦσιν, ἕτεροι δὲ νομέας ; Σ Od.~19.574 δρύοχοι δὲ ξύλα εἰσὶν ὀρθὰ ὑποκάτω τῆς τρόπιδος, ἐφ᾽ ὧν ἐπερείδεται, ἵνα μὴ αὐτὴν ἡ ψάμμος ἐσθίῃ. τινὲς δὲ δρυόχους φασὶ τὰ πρῶτα πηγνύμενα ξύλα εἰς ναυπηγίαν ; Eust.~1879 (ad Od.~19.574) Δρύοχοι δὲ κυρίως πάσσαλοι, ἐφ᾽ ὧν στοιχηδὸν διατεθειμένων ἡ τρόπις ἵσταται τῶν καινουργουμένων πλοίων διὰ ἰσότητα ; Hsch.~δ~2412 δρύακες· ⟨ἐπὶ⟩ τῶν ξύλων τῶν βασταζόντων τὴν τρόπιν τοῦ πλοίου ; Hsch.~ν~628 νομέας· ξύλα περιφερῆ ; Hsch.~ν~629 νομεῖς· βασιλεῖς, ἡγεμόνες, καὶ τῶν πλοίων τὰ ἐγκοίλια (Hdt. 1, 194, 2). καὶ σχοῖνοι ἀρμένων ; Hsch.~σ~1633 σταμίνες· παραστάται, καὶ τὰ ἐπὶ τῆς σχεδίας ὀρθὰ ξύλα, πρὸς ἃ αἱ σανίδες προσηλοῦνται, ἢ πάσσαλοι. παρὰ τὸ ἑστάναι (ε 252) ; Hsch.~τ~1498 τρόπις· τὸ κατώτατον τῆς νεώς (η 252 . .) ; EM~288.37 Δρυόχους, τοὺς πελέκεις. οἱ μὲν, κρίκους ἀκούουσι τινὰς μεγάλους ἐπ᾽ ὀβελίσκων κειμένους … ἄμεινον δὲ ἀκούειν δρυόχους ξύλα ὀρθά, ἐφ᾽ ὧν ἡ τρόπις ἐρείδεται τῆς πηγνυμένης νεώς ; Hsch.~δ~2432 δρύοχοι· τῶν σιδηρῶν πελέκεων αἱ ὀπαί, εἰς ἃς τὸ στελεὸν ἐνείρεται (τ 574) ; Hsch.~σ~1713 στεῖρα· ἡ ἄτοκος (κ 522). καὶ τὸ ἐξέχον τῆς πρῴρας ξύλον κατὰ τὴν τρόπιν (Α 482) ; Zen.~2.40 Ἅλις δρυός: ἐπὶ τῶν ἐκ φαυλοτέρας διαίτης ἐρχομένων ἐπὶ βελτίονα εἴρηται ἡ παροιμία}
\lp{52 δρυόχους τιθέναι]}{Ar.~\textit{Vesp.}~30 λέγε νυν ἀνύσας τι τὴν τρόπιν τοῦ πράγματος ; Pl.~\textit{Lg.}~7.803a–b οἷον δή τις ναυπηγὸς τὴν τῆς ναυπηγίας ἀρχὴν καταβαλλόμενος τὰ τροπιδεῖα ὑπογράφεται τῶν πλοίων σχήματα ; Ar.~\textit{Eq.}~542–544 ἐρέτην χρῆναι πρῶτα γενέσθαι πρὶν πηδαλίοις ἐπιχειρεῖν, κᾆτ᾽ ἐντεῦθεν πρῳρατεῦσαι καὶ τοὺς ἀνέμους διαθρῆσαι ; Theoc.~7.45–48 ὥς μοι καὶ τέκτων μέγ᾽ ἀπέχθεται, ὅστις ἐρευνῇ ἶσον ὄρευς κορυφᾷ τελέσαι δόμον εὐρυμέδοντος, καὶ Μοισᾶν ὄρνιχες}
\lp{52 τιθέναι δράματος ἀρχάς]}{Ar.~\textit{Ran.}~1119–1120 καὶ μὴν ἐπ᾽ αὐτοὺς τοὺς προλόγους σου τρέψομαι, ὅπως τὸ πρῶτον τῆς τραγῳδίας μέρος ; Pi.~\textit{O.}~6.1–4 χρυσέας ὑποστάσαντες εὐτειχεῖ προθύρῳ θαλάμου κίονας, ὡς ὅτε θαητὸν μέγαρον πάξομεν· ἀρχομένου δ᾽ ἔργου πρόσωπον χρὴ θέμεν τηλαυγές ; Pi.~\textit{P.}~4.138 βάλλετο κρηπῖδα σοφῶν ἐπέων ; Pl.~\textit{R.}~2.377a οὐκοῦν οἶσθ᾽ ὅτι ἀρχὴ παντὸς ἔργου μέγιστον, ἄλλως τε δὴ καὶ νέῳ καὶ ἁπαλῷ ὁτῳοῦν; μάλιστα γὰρ δὴ τότε πλάττεται ; Hsch.~π~3610 προοίμιον· πρόλογος, ἀρχὴ παντὸς λόγου ; Pi.~\textit{P.}~7.1–4 κάλλιστον αἱ μεγαλοπόλιες Ἀθᾶναι προοίμιον Ἀλκμανιδᾶν εὐρυσθενεῖ γενεᾷ κρηπῖδ᾽ ἀοιδᾶν ἵπποισι βαλέσθαι ; Hsch.~α~4186 ἀναβάλλετο· προοιμιάζετο (α 155 v. l.?) ; Hsch.~δ~2311 δραματουργεῖν· πανουργεῖν. συντιθέναι. πράττειν ; Hld.~7.6 ὥσπερ εἰς ἀνταγώνισμα δράματος ἀρχὴν ἄλλου παρεισφέρουσα ; Hsch.~ο~291 οἴμη· φωνή. ὁδός. λόγος. ἱστορία. ᾠδή. ὅθεν προοίμιον τὸ πρὸ τῆς ᾠδῆς. καὶ κύκλος ; Hsch.~φ~898 φροίμια· προοίμια, προρρήσεις. Ἀττικοί (Eur. Hec. 181)}
\lp{52 τιθέναι δράματος ἀρχάς (ἀρχὴ ἥμισυ παντός)]}{Diogenian.~2.97 Ἀρχὴ ἥμισυ παντός: ἐπὶ τῶν εἴς τι προκοπτόντων. Παρῆκται δὲ ἐκ τῆς Ἡσιόδου ; \textit{idem} Luc.~\textit{Somn.}~3 (ἀρχὴ δέ τοι ἥμισυ παντός), Suid.~α~4097 (ἐπὶ τῶν τινος ἀρχομένων καὶ προκοπτόντων ; om. Παρῆκται … Ἡσιόδου), Tosi n°~59 ; App. Prov.~1.41 Ἀρχὴ δήπου παντὸς ἔργου χαλεπώτερον ἐστί ; Ar.~\textit{Vesp.}~846 ἀφ᾽ Ἑστίας ἀρχόμενος ἐπιτρίψω τινά ; Pl.~\textit{Euthphr.}~3a ἀφ᾽ ἑστίας ἄρχεσθαι ; Zen.~1.40 Ἀφ᾽ Ἑστίας ἀρχόμενοι: μετενήνεκται δὲ ἀπὸ τῶν περὶ τὰ ἱερὰ δρωμένων ; \textit{idem} App. Prov.~1.43 (Ἀφ᾽ ἑστίας· ἐπὶ τῶν ἐξ ἀρχῆς τι πραττόντων) ; Zen.~2.7 Ἀπὸ βαλβῖδος: οἷον ἀπ᾽ ἀρχῆς χρὴ ἀγαθὸν εἶναι, ἢ τεχνίτην, ἢ κοινῶς ἐπὶ πάντων ; Erasm.~\textit{Adag.}~1.2.39 Principium dimidium totius. Ἀρχὴ ἥμισυ παντός, i. e. Principium dimidium totius. Hoc adagio significatum est, maximam difficultatis partem in aggrediundo negocio sitam esse ; Pl.~\textit{Lg.}~6.753e ἀρχὴ γὰρ λέγεται μὲν ἥμισυ παντὸς ἐν ταῖς παροιμίαις ἔργου, καὶ τό γε καλῶς ἄρξασθαι πάντες ἐγκωμιάζομεν ἑκάστοτε ; Hes.~\textit{Op.}~40 νήπιοι, οὐδὲ ἴσασιν ὅσῳ πλέον ἥμισυ παντὸς}
\lp{52 δράματος]}{Hsch.~δ~2310 δράματα· ποιήματα (Greg. Naz. or. 21, 14 [35, 1096C M.]) … πρᾶξις ἢ λόγος περιέχων ὑπόθεσίν τινα ; Ar.~\textit{Ach.}~415 δός μοι ῥάκιόν τι τοῦ παλαιοῦ δράματος ; Ar.~\textit{Ach.}~470 φροῦδά μοι τὰ δράματα ; Ar.~\textit{Pax}~795 εἶχε τὸ δρᾶμα γαλῆν τῆς ἑσπέρας ἀπάγξαι ; Ar.~\textit{Thesm.}~149–151 χρὴ γὰρ ποιητὴν ἄνδρα πρὸς τὰ δράματα / ἃ δεῖ ποιεῖν πρὸς ταῦτα τοὺς τρόπους ἔχειν … αὐτίκα γυναικεῖ᾽ ἢν ποιῇ τις δράματα ; Ar.~\textit{Thesm.}~849 τῷ δῆτ᾽ ἂν αὐτὸν προσαγαγοίμην δράματι ; Hdt.~6.21.2 ποιήσαντι Φρυνίχῳ δρᾶμα Μιλήτου ἅλωσιν καὶ διδάξαντι ἐς δάκρυά τε ἔπεσε τὸ θέητρον ; Hsch.~λ~1216 λόγος· ἡ τοῦ δράματος ὑπόθεσις}
\lp{53 κάμπτει]}{Ar.~\textit{Nub.}~333 κυκλίων τε χορῶν ᾀσματοκάμπτας ἄνδρας μετεωροφένακας ; Ar.~\textit{Thesm.}~68 ὄντος κατακάμπτειν τὰς στροφὰς οὐ ῥᾴδιον ; Pherecr.~fr.~155~K.-A. ap. [Plu.]~\textit{De\,mus.}~30 Κινησίας δὲ μ᾽ ὁ κατάρατος Ἀττικός, ἐξαρμονίους καμπὰς ποιῶν ἐν ταῖς στροφαῖς ἀπολώλεχ᾽ οὕτως, ὥστε τῆς ποιήσεως τῶν διθυράμβων, καθάπερ ἐν ταῖς ἀσπίσιν, ἀριστέρ᾽ αὐτοῦ φαίνεται τὰ δεξιά ; Pherecr.~fr.~155~K.-A. ap. [Plu.] De mus.~30 (Phrynis) Φρῦνις δ᾽ ἴδιον στρόβιλον ἐμβαλών τινα κάμπτων με καὶ στρέφων ὅλην διέφθορεν ; Hsch.~κ~617 κάμπτειν· τὸ ἐν τῇ ᾠδῇ καμπὰς ποιεῖν ; Poll.~7.119 τὸ δὲ τῆς οἰκοδομικῆς ἔργον οἰκοδομεῖν … λίθους ἁρμόττειν συναρμόττειν, ἀκριβῶς συμφυῶς ἀπισοῦν, ἀπευθύνειν, ἀποτείνειν, κάμπτειν, περιάγειν ; Ar.~\textit{Nub.}~969 εἰ δέ τις αὐτῶν βωμολοχεύσαιτ᾽ ἢ κάμψειέν τινα καμπήν ; Pherecr.~fr.~155~K.-A. ap. [Plu.] De mus.~30 (Philoxène) ὥσπερ τε τὰς ῥαφάνους ὅλην καμπῶν με κατεμέστωσε ; Timoth. ap. Plu.~\textit{De\,laude\,ipsius}~539c νικᾷ Τιμόθεος Μιλήσιος τὸν Κάμωνος τὸν ἰωνοκάμπταν ; Hsch.~κ~609 κάμπειος δρόμος· δρόμοι τινὲς ἦσαν κάμπειοι οὐκ εὐθεῖς καὶ ἁπλοῖ, ἀλλὰ καμπὰς ἔχοντες ; Hsch.~κ~608 καμπαλέας· καμπύλας ; Hsch.~κ~619 καμπύλα· ἐπικαμπῆ, στρεβλά (Prov. 2, 15) ; Phot.~τ~396 τὸ σκαμβὸν ξύλον οὐδέποτε ὀρθόν· αὕτη δημώδης ἐστὶ καὶ φανερά· τάττεται δὲ καὶ κατὰ τῶν διάστροφα ξύλα ἢ πράγματα πειρωμένων κατευθύνειν, μηδὲν ὠφελούντων}
\lp{53 κάμπτει δὲ νέας ἁψῖδας]}{\textit{Il.}~4.485–486 τὴν μέν θ᾽ ἁρματοπηγὸς ἀνὴρ αἴθωνι σιδήρῳ ἐξέταμ᾽, ὄφρα ἴτυν κάμψῃ περικαλλέϊ δίφρῳ ; Hes.~\textit{Op.}~426 τρισπίθαμον δ᾽ ἄψιν τάμνειν δεκαδώρῳ ἀμάξῃ ; Hsch.~α~8958 ἁψῖδες· τὰ κύκλα τῶν τροχῶν. αἱ περιφέρειαι. ἢ καμάραι ; Hsch.~ι~1100 ἴτυς· ἡ ἐσχάτη ἁψίς, ἢ περιαγωγὴ τῆς ἀσπίδος, καὶ τοῦ τροχοΰ (Δ 486), καὶ παντὸς περιφεροῦς τὸ τελευταῖον μέρος ; Eur.~\textit{Ba.}~1066–1067 κυκλοῦτο δ᾽ ὥστε τόξον ἢ κυρτὸς τροχὸς τόρνῳ γραφόμενος περιφορὰν ἕλκει δρόμον ; Eur.~\textit{Hipp.}~1233 ἁψῖδα πέτρῳ προσβαλὼν ὀχήματος ; Pi.~\textit{P.}~10.64–65 ὅσπερ ἐμὰν ποιπνύων χάριν τόδ᾽ ἔζευξεν ἅρμα Πιερίδων τετράορον ; Theoc.~25.247–249 ὡς δ᾽ ὅτ᾽ ἂν ἁρματοπηγὸς ἀνὴρ πολέων ἴδρις ἔργων ὄρπηκας κάμπτῃσιν ἐρινεοῦ εὐκεάτοιο, θάλψας ἐν πυρὶ πρῶτον, ἐπαξονίῳ κύκλα δίφρῳ ; Σ Ar.~\textit{Thesm.}~53 ὡς ἀπὸ τῶν οἰκοδομούντων ; Σ Pi.~\textit{P.}~10.102 τόδ᾽ ἔζευξεν ἅρμα: ἀντὶ τοῦ ζευχθῆναι ἐποίησεν. ἀλληγορικῶς δὲ τὸ ποίημα ἅρμα Πιερίδων ; Synag. B α~2625 ἁψίς· οὕτως ἐκάλουν τὸν ξύλινον κύκλον τοῦ τροχοῦ· τὸ δὲ ἐπ᾽ αὐτοῦ ἐπίσωτρον. Ὅμηρος δὲ τὸν μὲν τροχὸν κύκλον ὀνομάζει (Ψ 340), τὴν δὲ ἁψῖδα ἴτυν (e. g. Δ 486) ; Plu.~\textit{Cons.\,ad\,Apoll.}~103f τροχοῦ γὰρ περιστείχοντος ἄλλοθ᾽ ἡτέρα ἁψὶς ὕπερθε γίγνετ᾽ ἄλλοθ᾽ ἡτέρα ; \textit{idem} Macar.~8.58 (Τροχοῦ περιστείχοντος ἄλλοτε ἡ ἑτέρα ἁψὶς ἄνωθεν ἐστίν, ἄλλοτε ἡ ἑτέρα), Tosi n°~30 ; A.R.~3.137–138 χρύσεα μέν οἱ κύκλα τετεύχαται· ἀμφὶ δ᾽ ἑκάστῳ — διπλόαι ἁψῖδες περιηγέες εἱλίσσονται ; Hsch.~α~5545 ἄντυγος· περιφερείας τοῦ ἅρματος (Ε 262 . .) ; Hsch.~α~8968 ἁψῖσι· συναφαῖς (Ε 487) ; Hsch.~α~8969 ἁψῖσι λίνου· ἅμμασιν ἀπὸ τῆς συναφῆς (Ε 487) ; Hsch.~ο~1025 †ὀπίσσωτρον· ἡ ἀψὶς τοῦ τροχοῦ}
\lp{53 νέας … ἐπῶν]}{Ar.~\textit{Nub.}~1397 σὸν ἔργον ὦ καινῶν ἐπῶν κινητὰ καὶ μοχλευτὰ ; Eur.~\textit{Tro.}~512–514 ἀμφί μοι Ἴλιον, ὦ Μοῦσα, καινῶν ὕμνων ἄεισον ἐν δακρύοις ᾠδὰν ἐπικήδειον ; \textit{Od.}~1.351–352 τὴν γὰρ ἀοιδὴν μᾶλλον ἐπικλείουσ᾽ ἄνθρωποι, ἥ τις ἀκουόντεσσι νεωτάτη ἀμφιπέληται ; Timoth. ap. Ath.~3.122c οὐκ ἀείδω τὰ παλαιά ; Ar.~\textit{Ran.}~1381 τοὔπος νῦν λέγετον ἐς τὸν σταθμόν ; Pi.~\textit{O.}~9.48–49 αἴνει δὲ παλαιὸν μὲν οἶνον, ἄνθεα δ᾽ ὕμνων}
\lp{54 τὰ δὲ τορνεύει]}{\textit{Test.} Hsch.~τ~1168 τορνεύει· γλύφει ; Ar.~\textit{Av.}~490–491 χαλκῆς κεραμῆς σκυλοδέψαι σκυτῆς βαλανῆς ἀλφιταμοιβοὶ τορνευτολυρασπιδοπηγοί ; Ar.~\textit{Thesm.}~986 τόρευε πᾶσαν ᾠδήν ; \textit{Od.}~5.249 ὅσσον τίς τ᾽ ἔδαφος νηὸς τορνώσεται ἀνὴρ ; Hor.~\textit{Ars}~441 et male tornatos incudi reddere versus ; Pl.~\textit{Phdr.}~234e ὅτι σαφῆ καὶ στρογγύλα, καὶ ἀκριβῶς ἕκαστα τῶν ὀνομάτων ἀποτετόρνευται ; Hsch.~τ~1170 τόρνος· ἐργαλεῖον τεκτονικόν, ᾧ τὰ στρογγύλα σχήματα περιγράφεται ; Suid.~τ~789 Τορεύει: γλύφει, τιτρᾷ. καί, τόρευε πᾶσαν ᾠδήν. ἀντὶ τοῦ τορῶς καὶ τρανῶς λέγε τὴν ᾠδήν ; Ar.~\textit{Ran.}~818–821 σχινδαλάμων τε παραξόνια σμιλεύματά τ᾽ ἔργων, φωτὸς ἀμυνομένου φρενοτέκτονος ἀνδρὸς ῥήμαθ᾽ ἱπποβάμονα ; Eur.~\textit{Cyc.}~661 τόρνευ᾽, ἕλκε, μή σ᾽ ἐξοδυνηθεὶς δράσῃ τι μάταιον ; Thgn.~805–806 Τόρνου καὶ στάθμης καὶ γνώμονος ἄνδρα θεωρόν εὐθύτερόν με χρή, Κύρνε, φυλασσέμεναι ; Prop.~2.34.43 incipe iam angusto versus includere torno ; D.H.~\textit{Dem.}~43 αἳ μέν εἰσιν εὐκόρυφοι καὶ στρογγύλαι ὥσπερ ἀπὸ τόρνου ; Philostr.~\textit{VS}~1.25.7 καὶ ὅτε ἀποτορνεύοι περίοδον ; Jul.~\textit{Or.}~2.77a ἐργάτης γάρ ἐστι καὶ τούτων ἀγαθός, οὐκ ἀποσμιλεύων οὐδὲ ἀπονυχίζων τὰ ῥήματα οὐδὲ ἀποτορνεύων τὰς περιόδους καθάπερ οἱ κομψοὶ ῥήτορες ; Jul.~\textit{Or.}~3.112a λέξει τε ἡδίστῃ κοσμῶν καὶ ἐπιλεαίνων τὸ φαῦλον καὶ ἀγεννὲς τῶν διηγημάτων, γελοιότερον νομίζει τῶν ἀποτορνεύειν τὰς κέγχρους ἐπιχειρούντων, καθάπερ οἶμαι φασὶ τὸν Μυρμηκίδην ἀντιταττόμενον τῇ Φειδίου τέχνῃ ; \textit{idem} Suid.~γ~118 (καὶ παροιμία· γελοιότερον ἀπεργάζῃ ; ἀποτορεύειν ; ἢ καθάπερ ; ἀντιπραττόμενον) ; Hsch.~τ~1169 τορνευταί· γλύπται ; Hsch.~τ~1171 τορνῶσαι· περιγράψαι. κυκλῶσαι ; Hsch.~τ~1173 τορνοῦμαι δὲ πρὸς μέτρον· ἀντὶ τοῦ περιγράφομαι ; Phot.~τ~391 τόρνος· ξύλον στρογγύλον ; Erasm.~\textit{Adag.}~1.5.92 Incudi reddere. Eleganter Horatius, incudi reddere dixit, pro eo, quod est, refingere ac mutare, corrigereque: Et male tornatos incudi reddere versus ; Eur.~\textit{HF}~978 τόρνευμα δεινὸν ποδός ; Eur.~fr.~382 ap. Ath.~10.454b κύκλος τις ὡς τόρνοισιν ἐκμετρούμενος ; D.H.~\textit{Thuc.}~24 καθ᾽ ἓν ἕκαστον τῶν τῆς φράσεως μορίων ῥινῶν καὶ τορεύων ; Hsch.~τ~1172 τορνώσας· τῷ κύκλῳ περιγράψας (Nonn. Dion. 5, 63) ; Erasm.~\textit{Adag.}~2.10.16 Milium torno sculpere. Τὰς κέγχρους ἀποτορνεύειν. i. e. Milii grana torno sculpere dicebantur, qui rem absurdam moliuntur, Milium enim minutius est, quam ut inde quicquam desculpi possit}
\lp{54 τὰ δὲ κολλομελεῖ]}{Pi.~\textit{N.}~7.77–79 εἴρειν στεφάνους ἐλαφρόν· ἀναβάλεο· Μοῖσά τοι κολλᾷ χρυσὸν ἔν τε λευκὸν ἐλέφανθ᾽ ἁμᾷ ; Pl.~\textit{Menex.}~236b συνετίθει τὸν ἐπιτάφιον λόγον ὃν Περικλῆς εἶπεν, περιλείμματ᾽ ἄττα ἐξ ἐκείνου συγκολλῶσα ; Pl.~\textit{Phdr.}~278d–e ἄνω κάτω στρέφων ἐν χρόνῳ, πρὸς ἄλληλα κολλῶν τε καὶ ἀφαιρῶν ; Σ Ar.~\textit{Thesm.}~54 ὡς ἐπὶ τεκτόνων. μέλη κολλᾷ ; Hsch.~γ~616 Γλαύκου τέχνη· παροιμία ἐπὶ τῶν ῥαδίως κατεργαζομένων· Διονυσόδωρος δὲ τὴν περὶ τὸν σίδηρον κόλλησιν· Γλαῦκος γὰρ Χῖος σιδήρου κόλλησιν εὗρεν (Plat. Phaed. 108d) ; Plu.~\textit{De\,glor.\,Ath.}~350e ἀντίθετα καὶ πάρισα καὶ ὁμοιόπτωτα κολλῶν καὶ συντιθεὶς, μονονοὺ κολαπτῆρσι καὶ ξυστῆρσι τὰς περιόδους ἀπολεαίνων ; Hsch.~κ~3334 †κόλλα· τοῦ βοὸς τὸ νωτιαῖον δέρμα, ἐξ οὗ τὸ κολλᾶν ; Hsch.~μ~684 μέλη· τὰ γυῖα. καὶ ᾠδή. καὶ σκεύη ; Diogenian.~4.8 Γλαύκου τέχνη: ἐπὶ τῶν μὴ ῥᾳδίως κατεργαζομένων. Ἀπὸ Γλαύκου Σαμίου, ὃς πρῶτον κόλλησιν ἐφεῦρε σιδήρου ; Hdt.~1.25.2 ὑποκρητηρίδιον σιδήρεον κολλητόν, θέης ἄξιον διὰ πάντων τῶν ἐν Δελφοῖσι ἀναθημάτων, Γλαύκου τοῦ Χίου ποίημα, ὃς μοῦνος δὴ πάντων ἀνθρώπων σιδήρου κόλλησιν ἐξεῦρε ; Hsch.~κ~3336 κολλήεντα· †κακῶς κεκολλημένα, ἢ τὰ ἐκ συμβλημάτων ἡρμοσμένα (Ο 389) ; Hsch.~κ~3337 κολλητὰς σανίδας· τὰς θύρας, τὰ θυρώματα (Ι 583) ; Hsch.~κ~3344 κολλοπῶσαι· κατακολλῆσαι ; Erasm.~\textit{Adag.}~2.8.30 Glauci ars. Γλαύκου τέχνη. i. e. Glauci ars. Hujus paroemiae duplex est usus. Convenit enim vel in opus celeriter absolutum, vel in id, quod est arte summa, summaque cura elaboratum}
\lp{55 γνωμοτυπεῖ κἀντονομάζει]}{Ar.~\textit{Nub.}~321 καὶ γνωμιδίῳ γνώμην νύξασ᾽ ἑτέρῳ λόγῳ ἀντιλογῆσαι ; Agatho ap. Arist. EN~6.4 (1140a19) καθάπερ καὶ Ἀγάθων φησὶ τέχνη τύχην ἔστερξε καὶ τύχη τέχνην ; Agatho ap. Ath.~5.185a κατὰ τὸν καλὸν Ἀγάθωνα· τὸ μὲν πάρεργον ἔργον ὣς ποιούμεθα, τὸ δ᾽ ἔργον ὡς πάρεργον ἐκπονούμεθα ; Agatho ap. Arist.~\textit{Rh.}~2.24.10 ὥσπερ καὶ Ἀγάθων λέγει τάχ᾽ ἄν τις εἰκὸς αὐτὸ τοῦτ᾽ εἶναι λέγοι, βροτοῖσι πολλὰ τυγχάνειν οὐκ εἰκότα ; Agatho ap. Ath.~13.584a κατὰ γὰρ τὸν Ἀγάθωνα· γυνὴ τοι σώματος δι᾽ ἀργίαν ψυχῆς φρόνησιν ἐντὸς οὐκ ἀργὸν φορεῖ}
\lp{55 γνωμοτυπεῖ]}{Ar.~\textit{Eq.}~1379 καὶ γνωμοτυπικὸς καὶ σαφὴς καὶ κρουστικός ; Ar.~\textit{Nub.}~950 λόγοισι καὶ φροντίσι καὶ γνωμοτύποις μερίμναις ; Ar.~\textit{Ran.}~876–877 ἀνδρῶν γνωμοτύπων, ὅταν εἰς ἔριν ὀξυμερίμνοις ; Arist.~\textit{Rh.}~2.21.9 σημεῖον δὲ ἱκανόν· οἱ γὰρ ἀγροῖκοι μάλιστα γνωμοτύποι εἰσὶ καὶ ῥᾳδίως ἀποφαίνονται ; Cratin.~fr.~342~K.-A. (=~307~Kock) τίς δὲ σύ; κομψός τις ἔροιτο θεατής. ὑπολεπτολόγος, γνωμοδιώκτης, εὐριπιδαριστοφανίζων ; Aesch.~\textit{Ag.}~1153 τὰ δ᾽ ἐπίφοβα δυσφάτῳ κλαγγᾷ μελοτυπεῖς ὁμοῦ τ᾽ ὀρθίοις ἐν νόμοις ; Pi.~\textit{P.}~1.86 ἀψευδεῖ δὲ πρὸς ἄκμονι χάλκευε γλῶσσαν ; Pl.~\textit{Phdr.}~267c τὰ δὲ Πώλου πῶς φράσωμεν αὖ μουσεῖα λόγων—ὡς διπλασιολογίαν καὶ γνωμολογίαν καὶ εἰκονολογίαν ; Arist.~\textit{Rh.}~3.19.1 (1419b) πέφυκε γάρ, μετὰ τὸ ἀποδεῖξαι αὐτὸν μὲν ἀληθῆ τὸν δὲ ἐναντίον ψευδῆ, οὕτω τὸ ἐπαινεῖν καὶ ψέγειν καὶ ἐπιχαλκεύειν ; Plu.~\textit{De\,garr.}~17,~511b θαυμάσαντες τῆς λέξεως τὸ εὔογκον καὶ τὸ λιτόν, ἐν βραχεῖ σφυρήλατον νοῦν περιεχούσης ; Hsch.~γ~747 γνωμοτυπία· γνωμολογία ; Suid.~γ~345 Γνωμοτυπικῶς: γνωμολογικῶς ; Ar.~\textit{Eq.}~100 βουλευματίων καὶ γνωμιδίων καὶ νοιδίων ; Plu.~\textit{De\,Pyth.\,or.}~29,~408e οὐχ ἣκιστα διὰ τὴν βραχυλογίαν ὡς πυκνὸν καὶ σφυρήλατον νοῦν ἐν ὀλίγῳ περιέχουσαν ὄγκῳ ; [Luc.]~\textit{Dem.\,enc.}~14 μόνος γέ τοι τῶν ῥητόρων, ὡς ὁ Λεωσθένης ἐτόλμησεν εἰπεῖν, ἔμψυχον καὶ σφυρήλατον παρεῖχε τὸν λόγον ; Hsch.~γ~745 γνῶμαι· τὰ ψηφίσματα q ; App. Prov.~1.81 Γνῶθι σαυτόν καὶ Τὸ νόμισμα παραχάραξον: παραγγέλματα Πυθικά}
\lp{55 κἀντονομάζει]}{\textit{Test.} Suid.~α~2754 Ἀντονομάζει: ὀνόματα ἀντὶ ὀνομάτων τίθησι ; \textit{idem} Σ Ar.~\textit{Thesm.}~55 (ὄνομα ἀντὶ ὀνόματος ; om. Ἀντονομάζει) ; Ar.~fr.~341~K.-A. καὶ κατ᾽ Ἀγάθων᾽ ἀντίθετον ἐξευρημένον ; Ar.~fr.~599~Kock ap. Ath.~15.701b ἐκφέρετε πεύκας κατ᾽ Ἀγάθωνα φωσφόρους ; Ael.~\textit{VH}~14.13 πολλοῖς καὶ πολλάκις χρῆται τοῖς ἀντιθέτοις ὁ Ἀγάθων. ἐπεὶ δέ τις οἷον ἐπανορθούμενος αὐτὸν ἐβούλετο περιαιρεῖν αὐτὰ τῶν ἐκείνου δραμάτων, εἶπεν ʽἀλλὰ σύ γε, γενναῖε, λέληθας σεαυτὸν τὸν Ἀγάθωνα ἐκ τοῦ Ἀγάθωνος ἀφανίζων.᾽ ; Ath.~5.187c χλευάζει τε τὰ ἰσόκωλα τὰ Ἀγάθωνος καὶ τὰ ἀντίθετα ; Th.~6.4.6 τὴν πόλιν αὐτὸς ξυμμείκτων ἀνθρώπων οἰκίσας Μεσσήνην ἀπὸ τῆς ἑαυτοῦ τὸ ἀρχαῖον πατρίδος ἀντωνόμασεν ; Pl.~\textit{Cra.}~409c συγκεκροτημένον δὲ Σελαναία κέκληται. … διθυραμβῶδές γε τοῦτο τοὔνομα, ὦ Σώκρατες ; Hermog.~\textit{Stat.}~4 (p.~62.1–3 Rabe) καὶ πρῶτός γέ ἐστιν ὁ ἀντονομάζων καλούμενος, ὅταν ὄνομα ὀνόματι ὁ φεύγων ἀντιτιθῇ ; Tryph.~\textit{Trop.}~31 Ἀντονομασία ἐστὶ λέξις ἢ φράσις διὰ συνωνύμων ὀνομάτων τὸ κύριον παριστῶσα· λέξις μὲν Φοῖβε ἀντὶ τοῦ Ἄπολλον ; Hsch.~α~5413 ἀντίθετα· ἐναντία}
\lp{56–57 κηροχυτεῖ … χοανεύει]}{Ar.~fr.~719~K.-A. (=~699~Kock) ῥήματά τε κομψὰ καὶ παίγνι᾽ ἐπιδεικνύναι πάντ᾽ ἀπ᾽ ἀκροφυσίων κἀπὸ καναβευμάτων ; Poll.~10.189 αὐτὸ δὲ τὸ πήλινον, ὃ περιείληφε τὰ πλασθέντα κήρινα, ἃ κατὰ τὴν τοῦ πυρὸς προσφορὰν τήκεται καὶ πολλὰ ἐκείνῳ τρυπήματα ἐναπολείπεται, μίλιγδος καλεῖται ; Suid.~α~2874 λέγει γὰρ διὰ μὲν τοῦ ἀπ᾽ ἀκροφυσίων, καινῶς εἰργασμένα καὶ οἷον ἐκ πυρός, διὰ δὲ τοῦ ἀπὸ κινναβευμάτων, οἷον καινῶς πεπλασμένα καὶ διάθεσιν ἔχοντα. κίνναβος γὰρ τὸ εἴδωλον, πρὸς ὃ οἱ πλάσται καὶ οἱ ζωγράφοι βλέποντες διατίθενται πλάττοντες καὶ γράφοντες ; Plu.~\textit{De\,superst.}~6,~167e εἶτα χαλκοτύποις μὲν πείθονται καὶ λιθοξόοις καὶ κηροπλάσταις ἀνθρωπόμορφα τῶν θεῶν τὰ εἴδη ποιοῦσι}
\lp{56 καὶ κηροχυτεῖ]}{Ar.~\textit{Nub.}~149–150 κηρὸν διατήξας, εἶτα τὴν ψύλλαν λαβὼν ἐνέβαψεν ἐς τὸν κηρὸν αὐτῆς τὼ πόδε ; Castorion ap. Ath.~10.454f–455b μουσοπόλε θήρ, κηρόχυτον ὃς μείλιγμ᾽ ἱεῖς ; Pl.~\textit{R.}~9.588d δεινοῦ πλάστου, ἔφη, τὸ ἔργον· ὅμως δέ, ἐπειδὴ εὐπλαστότερον κηροῦ καὶ τῶν τοιούτων λόγος, πεπλάσθω ; Hippod. ap. Stob.~4.1.94 τὰ δὲ ἔθεα καὶ ἐπιταδεύματα πλάσσει καὶ κηροχυτεῖ τὰν ψυχάν ; Cic.~\textit{De\,or.}~3.177 sicut mollissimam ceram ad nostrum arbitrium formamus et fingimus ; Suid.~κ~1537 Κηροῦ εὐπλαστότερος: ἐπὶ τῶν ῥᾳδίως τι ποιούντων ; \textit{idem} Greg. Cypr.~\textit{L}~2.44, Tosi n°~1375 ; Ar.~\textit{Ran.}~1040 ὅθεν ἡμὴ φρὴν ἀπομαξαμένη πολλὰς ἀρετὰς ἐποίησεν ; Aesch.~\textit{PV}~574 ὑπὸ δὲ κηρόπλαστος ὀτοβεῖ δόναξ ; \textit{AP}~16.210 ξουθαὶ δ᾽ ἐφύπερθε μέλισσαι κηροχύτου μέλιτος ; Pl.~\textit{Tht.}~191c θὲς δή μοι λόγου ἕνεκα ἐν ταῖς ψυχαῖς ἡμῶν ἐνὸν κήρινον ἐκμαγεῖον ; Pl.~\textit{Ti.}~74d ταῦτα ἡμῶν διανοηθεὶς ὁ κηροπλάστης, ὕδατι μὲν καὶ πυρὶ καὶ γῇ συμμείξας καὶ συναρμόσας, ἐξ ὀξέος καὶ ἁλμυροῦ συνθεὶς ζύμωμα ὑπομείξας αὐτοῖς, σάρκα ἔγχυμον καὶ μαλακὴν συνέστησεν ; Hsch.~κ~628 κανάβιος κηρός· ᾧ χρῶνται οἱ ἀνδριαντοποιοὶ πρὸς πλάσιν ; Hsch.~κ~629 κάναβοι· τὰ ξύλα, περὶ ἃ τὸ πρῶτον οἱ πλάσται τὸν κηρὸν τιθέασιν. ὅθεν καὶ οἱ λεπτοὶ καὶ ἄσαρκοι κάναβοι λέγονται ; Hsch.~κ~2546 κηρίον· τὸ τῶν μελισσῶν. καὶ εἶδος πλακοῦντος ; Poll.~7.163–165 καὶ κηροπλάσται κηροπλαστική, κηρὸν τήκειν ; Erasm.~\textit{Adag.}~3.2.33 Cera tractabilior. … De vehementer docili dicebatur, aut de ingenio facili, tractabilique, \& in quemvis habitum sequaci. Horatius: Cereus in vitium flecti, monitoribus asper. Persius quinta Satyra: Artificemque tuo ducit sub pollice vultum. … Est autem hoc cerae peculiare, ut manibus contrectata mollescat. Et hinc variae rerum formae solent effingi ; Eub.~\textit{vel\,Araros} ap. Ath.~13.562d τίς ἦν ὁ γράψας πρῶτος ἀνθρώπων ἄρα ἢ κηροπλαστήσας Ἔρωθ᾽ ὑπόπτερον ; Phot. s.v.~κάναβος (I, p.~319 Naber) κάναβος· τὸ πρῶτον τιθέμενον ξύλον ὑπὸ τῶν πλαστῶν, περὶ ὃ τὸν πηλὸν τιθέντες πλάσσουσιν}
\lp{56 καὶ γογγύλλει]}{\textit{Test.} Hsch.~γ~764 γογγύλλειν· συστρέφειν (Ar. Thesm. 56) ; Suid.~γ~361 Γογγύλιον: γύλιον. στρογγύλον. καὶ Γογγυλίζει, ἀντὶ τοῦ μεταστρέφει ; \textit{idem} Zonar. col.~449~Tittmann (Γογγυλίζει. ἀντὶ τοῦ μεταστρέφει tantum) ; Ar.~\textit{Ach.}~686 ἐς τάχος παίει ξυνάπτων στρογγύλοις τοῖς ῥήμασιν ; Ar.~\textit{Pax}~28 ὥσπερ γυναικὶ γογγύλην μεμαγμένην ; Ar.~\textit{Lys.}~973–975 εἴθ᾽ αὐτὴν ὥσπερ τοὺς θωμοὺς μεγάλῳ τυφῷ καὶ πρηστῆρι ξυστρέψας καὶ ξυγγογγύλας ; Ar.~fr.~488~K.-A. (=~471~Kock) χρῶμαι γὰρ αὐτοῦ τοῦ στόματος τῷ στρογγύλῳ, τοὺς νοῦς δ᾽ ἀγοραίους ἧττον ἢ ᾽κεῖνος ποιῶ ; Demetr.~\textit{Eloc.}~20 Τῆς δὲ ῥητορικῆς περιόδου συνεστραμμένον τὸ εἶδος καὶ κυκλικὸν καὶ δεόμενον στρογγύλου στόματος καὶ χειρὸς συμπεριαγομένης τῷ ῥυθμῷ ; Σ Ar.~\textit{Thesm.}~56 ὡς ἐπιπηγνύντος καὶ συντιθέντος. μεταστρέφει ; Hsch.~γ~763 γογγύλον· στρογγύλον. σκληρόν (Soph. fr. 395, 2 R.) ; EM~237.49 Γογγυλίζω, μεταστρέφω, παρὰ τὸ γογγύλη, ἥ ἐστιν εἶδος λαχάνου ; Ar.~\textit{Thesm.}~1185 οἴμ᾽ ὠς στέριπο τὸ τιττί, ὤσπερ γογγύλη ; Alex. ap. Ath.~4.165b–e σφαῖραν ἐποίησε τὴν πατρῴαν οὐσίαν οὕτως συνεστρόγγυλεν ἰταμῶς καὶ ταχύ ; Arist.~\textit{Rh.}~2.21.7 περὶ δὲ τῶν μὴ παραδόξων ἀδήλων δὲ προστιθέντα τὸ διότι στρογγυλώτατα ; D.H.~\textit{Dem.}~18 ἡ δ᾽ ἐναγώνιος στρογγύλη τε εἶναι βούλεται καὶ συγκεκροτημένη καὶ μηδὲν ἔχουσα κολπῶδες ; Hsch.~γ~759 γογγρύζειν· τονθρύζειν. τὸ ὡς ὓς φωνεῖν, ὅπερ ἔνιοι γογγύζειν. τὸ αὐτὸ καὶ γρυλίζειν λέγεται ; Hsch.~γ~765 γογγυλόσκηνος· στρογγύλον ἔχων τὸν οἶκον, ἢ τὸ σῶμα ; Hsch.~γ~767 γογγυλωπά· στρογγύλα, εὔτονα. γοργά ; Hsch.~σ~1992 στρέφει· κάμπτει (Ψ 323). λυγίζει ; Hsch.~σ~2020 στρογγύλη· συνεστραμμένη. καὶ Λιπάρα νῆσος ; Phryn.~\textit{PS}~59 de Borries γογγυλίς: οὕτως οἱ Ἀττικοὶ λέγουσιν, ἀλλ᾽ οὐ γογγύλην. γογγύλην δὲ τὸ στρογγύλον ξύλον ; Prov. Aesop.~11 (CPG II, p.~229) Στρογγύλα λέγε, ἵνα καὶ κυλίηται ; Nicom. ap. Ath.~2.58a οὐσίδιον μοι καταλιπόντος τοῦ πατρός, οὕτω συνεστρόγγυλα κἀξεκόκκισα ; Soph.~\textit{Ichn.}~305 πλεῖστον με[τ]αξύ· γογγύλον γάρ ἐστι καὶ βραχυσκελές ; Hsch.~γ~766 γογγυλεύματα· στρογγυλεύματα}
\lp{57 καὶ χοανεύει]}{\textit{Test.} Eust.~1153.40–41 (ad Il.~18.470) ἀπὸ δὲ τούτου καὶ παρὰ τῷ κωμικῷ τὸ χοανεύειν, ῥηθὲν οὐ σεμνῶς ἐφ᾽ οὗ ἐκεῖνος ἤθελε ; Hes.~\textit{Th.}~862–864 ἐτήκετο κασσίτερος ὣς τέχνῃ ὕπ᾽ αἰζηῶν ἐν ἐυτρήτοις χοάνοισι θαλφθείς ; Pl.~\textit{R.}~3.411a ὅταν μέν τις μουσικῇ παρέχῃ καταυλεῖν καὶ καταχεῖν τῆς ψυχῆς διὰ τῶν ὤτων ὥσπερ διὰ χώνης ἃς νυνδὴ ἡμεῖς ἐλέγομεν τὰς γλυκείας τε καὶ μαλακὰς καὶ θρηνώδεις ἁρμονίας ; Hsch.~γ~653 †γλίδου χοάναι· εἰς οὓς τύπους κατέρχεται τὸ ἐν τῇ χωνῃ διὰ τῶν κεντημάτων. ἦσαν δὲ πήλινοι. Σοφοκλῆς (fr. 35 R.) ; Hsch.~λ~960 λιγδαρεοχύται· οἱ ἐν τοῖς λίγδοις τὰς †σάρκας χέοντες, τουτέστιν χοάναις ; Hsch.~χ~576 χοάνη· χώνη, τύπος, εἰς ὃν μεταχεῖται τὸ χωνευόμενον (Ar. Thesm. 18) ; \textit{Il.}~18.470 φῦσαι δ᾽ ἐν χοάνοισιν ἐείκοσι πᾶσαι ἐφύσων ; Emp.~fr.~96.1–3~DK ap. Arist.~\textit{de\,An.}~1.5,~410a4–6 ἡ δὲ χθὼν ἐπίηρος ἐν εὐστέρνοις χοάνοισιν τὼ δύο τῶν ὀκτώ μερέων λάχε νήστιδος αἴγλης, τέσσαρα δ᾽ Ἡφαίστοιο· τὰ δ᾽ ὀστέα λευκὰ γένοντο ; Σ Ar.~\textit{Thesm.}~57 χωνεύει ; Σ Il.~18.470 χοάνοισιν] ἄδηλον τὸ γένος· λέγει δὲ ‘ τοῖς ἀγγείοις, εἰς ἃ τήκεται ὁ χαλκός,᾽ ἅπερ ἐστὶ παῤ ἡμῖν πήλινα ; Hsch.~λ~965 λίγδος· τόπος χοάνης, καὶ ἡ θυεία, καὶ ἐν ᾗ χωνεύουσιν (Soph. fr. 35 R.) ; Hsch.~λ~966 λίγδου χοάνη· οἱ δὲ τὰ λίκνα τῶν ἀργυρείων ; Hsch.~χ~577 χοάνοις· τοῖς φυσητῆρσι, ταῖς χώναις, καὶ κοιλώμασιν, εἰς ἃ ἐγχεῖται τὸ χωνευόμενον, ἢ τοῖς πηλίνοις τύποις ; A.R.~3.1299 ὡς δ᾽ ὅτ᾽ ἐνὶ τρητοῖσιν ἐύρρινοι χοάνοισιν ; Hp.~\textit{Cord.}~8 παρέθηκεν αὐτέῳ φύσας, καθάπερ τοῖσι χοάνοισιν οἱ χαλκέες, ὥστε διὰ τουτέων χειροῦται τὴν πνοήν ; Hsch.~σ~2937 σφυρήλατος· σφύραις ἐληλα[σ]μένος, οὐ χωνευτός ; Hsch.~σ~2938 σφυρήλατα· τὰ σιδήρια, ὅτι οὐ χεῖται ; Hsch.~χ~864 χωνευτήριον· χων⟨ε⟩ῖον (Zach. 11, 13) ; Synag.~χ~88 χοάνης· χωνίου}
\lp{57 καὶ λαικάζει]}{Ar.~\textit{Ach.}~79 ἡμεῖς δὲ λαικαστάς τε καὶ καταπύγονας ; Ar.~\textit{Eq.}~167 δήσεις φυλάξεις, ἐν πρυτανείῳ λαικάσει ; Ar.~\textit{Eq.}~1378–1381 συνερτικὸς γάρ ἐστι καὶ περαντικός, καὶ γνωμοτυπικὸς καὶ σαφὴς καὶ κρουστικός, καταληπτικός τ᾽ ἄριστα τοῦ θορυβητικοῦ. — οὔκουν καταδακτυλικὸς σὺ τοῦ λαλητικοῦ ; Ar.~\textit{Ran.}~1308 αὕτη ποθ᾽ ἡ Μοῦσ᾽ οὐκ ἐλεσβίαζεν, οὔ ; Strato ap. Ath.~9.382b–383b πηγὸς πάρεστι; πηγός; οὐχὶ λαικάσει … σαφέστερόν θ᾽ ὃ βούλει μοι λέγειν ; Hsch.~λ~692 λεσβιάζειν· πρὸς ἄνδρα στόμα μολύνειν· Λεσβιάδας γὰρ τὰς λαικαστρίας ἔλεγον ; Ar.~\textit{Ach.}~529 Ἕλλησι πᾶσιν ἐκ τριῶν λαικαστριῶν ; Ar.~\textit{Vesp.}~1346 μέλλουσαν ἤδη λεσβιᾶν τοὺς ξυμπότας ; Pherecr.~fr.~149~Kock (Chiron) ap. Σ Ar. Ran.~1308 δώσει δέ σοι γυναῖκας ἑπτὰ Λεσβίδας. — καλόν γε δῶρον, ἑπτ᾽ ἔχειν λαικαστρίας ; Cephisod. ap. Ath.~15.689e–690a ὦ λακκόπρωκτε, βάκκαριν τοῖς σοῖς ποσὶν ἐγὼ πρίωμαι; λαικάσομ᾽ ἄρα βάκκαριν ; Σ Ar.~\textit{Eccl.}~920 λάβδα : λαικάζουσιν οἱ Λέσβιοι ἀπὸ τοῦ ἄρχοντος στοιχείου ; Hsch.~λ~848 ληκᾶσθαι· περαίνεσθαι. καὶ τὰ αὐτά ; Hsch.~λ~858 ληκώ· τὸ μόριον ; Hsch.~σ~783 σιφνιάζειν· καταδακτυλίζειν· διαβέβληνται γὰρ οἱ Σίφνιοι ὡς παιδικοῖς χρώμενοι· σιφνιάσαι οὖν τὸ σκιμαλίσαι ; Hsch.~σ~929 σκερολίγγες· λαικασταί. ἢ †ὠπισταί† ; Ar.~\textit{Ach.}~537 Λακεδαιμονίων ἐδέοντο τὸ ψήφισμ᾽ ὅπως μεταστραφείη τὸ διὰ τὰς λαικαστρίας ; Ar.~\textit{Lys.}~17 ἡ μὲν γὰρ ἡμῶν περὶ τὸν ἄνδρ᾽ ἐκύπτασεν ; Luc.~\textit{Lex.}~12 μίνθων ἐκεῖνός γε καὶ λαικαλέος ; Suid.~λ~180 Λαικάζω: ἀπατῶ}
\lp{58 τίς ἀγροιώτας πελάθει θριγκοῖς;]}{Eur.~\textit{Ion}~154–157 φοιτῶσ᾽ ἤδη λείπουσίν τε πτανοὶ Παρνασοῦ κοίτας· αὐδῶ μὴ χρίμπτειν θριγκοῖς μηδ᾽ ἐς χρυσήρεις οἴκους ; Eur.~\textit{Ion}~170a–172 τίς ὅδ᾽ ὀρνίθων καινὸς προσέβα; μῶν ὑπὸ θριγκοὺς εὐναίας καρφυρὰς θήσων τέκνοις ; Eur.~\textit{Or.}~1268–1270 τίς ὅδ᾽ ἄρ᾽ ἀμ|φὶ μέλαθρον πολεῖ σὸν ἀγρότας ἀνήρ ; Sapph.~fr.~57 ap. Ath.~1.21c τίς δ᾽ ἀγροιῶτις θέλγει νόον οὐκ ἐπισταμένη τὰ βράκε᾽ ἕλκειν ἐπὶ τῶν σφυρῶν ; Eur.~\textit{IT}~128–129 πρὸς σὰν αὐλάν, εὐστύλων ναῶν χρυσήρεις θριγκούς ; [Eur.]~\textit{Rh.}~557–558 τί ποτ᾽ οὐ πλάθει σκοπός, ὃν ναῶν Ἕκτωρ ὤτρυνε κατόπταν}
\lp{58 ἀγροιώτας]}{Ar.~\textit{Thesm.}~159–160 ἄλλως τ᾽ ἄμουσόν ἐστι ποιητὴν ἰδεῖν ἀγρεῖον ὄντα καὶ δασύν ; \textit{Od.}~21.85 νήπιοι ἀγροιῶται, ἐφημέρια φρονέοντες ; Pl.~\textit{Smp.}~194c οὐ μεντἂν καλῶς ποιοίην, φάναι, ὦ Ἀγάθων, περὶ σοῦ τι ἐγὼ ἄγροικον δοξάζων ; Hsch.~α~820 ἀγροιῶται· ἄγροικοι (λ 293). καὶ γένος Ἀθήνῃσιν, ὃ ἀντιδιέστελλον πρὸς τοὺς Εὐπατρίδας ; Hsch.~α~821 †ἀγροιτιᾷ· ἀγροικεύεται. ὑβρίζει ; Suid.~α~380 Ἀγροιώτης: ὁ γεωργικὸς λεώς, ὁ ἐν ἀγρῷ. ἀγροιώτᾳ τῷδε μονοστόρθυγγι Πριήπῳ ; Zen.~1.15 Ἀγροίκου μὴ καταφρόνει ῥήτορος; ὅτι μηδὲ τῶν εὐτελῶν χρὴ καταφρονεῖν, παραινεῖ ἡ παροιμία ; Apostol.~1.24a Ἄγροικός εἰμι τὴν σκάφην σκάφην λέγων ; \textit{idem} Apostol.~15.95b (Τὰ σῦκα σῦκα, τὴν σκάφην σκάφην λέγει), Tosi n°~2272 (Τὰ σῦκα σῦκα, τὴν σκάφην δὲ σκάφην ὀνομάζων), Erasm.~\textit{Adag.}~2.3.5 (Τὰ σῦκα σῦκα, τὴν σκάφην σκάφην λέγων) ; Ar.~\textit{Nub.}~628 οὐκ εἶδον οὕτως ἄνδρ᾽ ἄγροικον οὐδένα ; Hes.~\textit{Sc.}~39 οὐδ᾽ ὅ γ᾽ ἐπὶ δμῶας καὶ ποιμένας ἀγροιώτας ; Σ Eur.~\textit{Or.}~1270 ἀγρότας ἀνήρ : ἀγροῖκος, δρυτόμος ; Hsch.~α~816 ἀγροῖκος· ὁ ἐν ἀγρῷ διάγων, χωρικός. ἢ ἐργάτης καὶ δραστήριος. ἢ ζευγηλάτης ; Hsch.~α~829 ἀγρότας· ἀγροίκους ; Hsch.~α~834 ἀγρόται· ἀγροῖκοι. ἢ θηρευταί ; Hsch.~α~844 ἀγρῶσται· ἐργάται (Eur. Rhes. 287). θηρευταί (Ap. Rh. 4, 175) ; Hsch.~α~847 ἀγρωστῖνος· ἀγροῖκος (Epicharm. tit., PCG 1, 17) ; Plu.~\textit{Mor.}~178b σκαιούς ἔφη φύσει καὶ ἀγροίκους εἶναι Μακεδόνας καὶ τὴν σκάφην σκάφην λέγοντας ; [Eur.]~\textit{Rh.}~287 φόβον δ᾽ ἀγρώσταις, οἳ κατ᾽ Ἰδαῖον λέπας ; \textit{Il.}~11.549 ἐσσεύαντο κύνες τε καὶ ἀνέρες ἀγροιῶται ; \textit{Od.}~11.293 δεσμοί τ᾽ ἀργαλέοι καὶ βουκόλοι ἀγροιῶται ; Theoc.~13.44 ἀκοίμητοι, δειναὶ θεαὶ ἀγροιώταις ; Theoc.~25.23 εὐθὺς δὲ σταθμοὶ περιμήκεες ἀγροιώταις ; A.R.~4.1183 σὺν δ᾽ ἀνέρες ἀγροιῶται ; Hsch.~α~753 ἄγραυλον· ὕπαιθρον, καὶ ἔρημον. ἢ ἐν ἀγρῷ αὐλιζόμενον ἢ καπυρόν ; Hsch.~α~817 ἀγροικότατον· ἀπαιδευτότατον ; Suid.~α~375 Ἄγροικος, ἄφρων, δύσκολος· σκληρὸς καὶ ἀπαίδευτος, ἢ ὁ ἐν ἀγρῷ κατοικῶν}
\lp{58 πελάθει]}{Aesch.~fr.~132.1–2~Radt ap. Ar.~\textit{Ran.}~1264–1265 Φθιῶτ᾽ Ἀχιλλεῦ, τί ποτ᾽ ἀνδροδάικτον ἀκούων ἰὴ κόπον οὐ πελάθεις ἐπ᾽ ἀρωγάν ; Hsch.~π~2433 πλάθεται· παραγίνεται, προσπελάζει (Preger 150, 2 = FGE 1669) ; Eur.~\textit{El.}~1292–1293 θέμις ἐς φθογγὰς τὰς ὑμετέρας ἡμῖν πελάθειν ; Hsch.~π~2431 πλαθείς· προσπελασθείς. δοθείς ; Hsch.~π~2432 πλαθεῖσα· τὰ αὐτά (Eur. e. g. Tr. 203) ; Hsch.~π~1278 πελάζειν· ἐγγίζειν, πλησιάζειν (Ε 766) ; Hsch.~π~1291 πελάσαι· ἐγγίσαι, πλησιάσαι (χ 176) ; Hsch.~π~1294 πέλασε· προσήγγισεν (Δ 123 . .) ; Synag.~π~272 πελάζει· πλησιάζει, ἐγγίζει}
\lp{58 θριγκοῖς]}{\textit{Test.} Hsch.~θ~750 θριγκοῖς· φραγμοῖς ; \textit{idem} Σ Ar.~\textit{Thesm.}~58 (τοῖς φραγμοῖς τῶν οἴκων ; om. θριγκοῖς) ; Eur.~\textit{Hel.}~430 ἰδὼν δὲ δῶμα περιφερὲς θριγκοῖς τόδε ; Hsch.~θ~749 θριγκός· ἡ στεφάνη τοῦ τείχους. περίφραγμα καὶ περίβολος (η 87). ἤγουν τὸ ἀνώτατον ⟨τῆς⟩ τοῦ τείχους οἰκοδομῆς ἐφ᾽ οὗ καὶ ἡ στέγη κεῖται. καὶ ὁ ὑπὸ τὸν κλινόποδα τοῦ τοίχου τόπος. ἢ περιχαράκωμα. καὶ τριγχός ; Eur.~\textit{El.}~1150 ἰάχησε δὲ στέγα λάινοί τε θριγκοὶ δόμων ; Eur.~\textit{IT}~74 θριγκοῖς δ᾽ ὑπ᾽ αὐτοῖς σκῦλ᾽ ὁρᾷς ἠρτημένα ; \textit{Od.}~7.86–87 ἐς μυχὸν ἐξ οὐδοῦ, περὶ δὲ θριγκὸς κυάνοιο ; Hsch.~γ~272 γεῖσον· στεφάνωμα οἴκου (Eur. Or. 1570) ; Hsch.~τ~369 τειχίον· Ἀττικοὶ τὸν περίβολον τοῖς χωρίοις ; Hsch.~τ~1365 τριγκῶσαι· φράξαι ; Suid.~θ~497 Θριγκός: τὸ περίφραγμα τοῦ οἴκου … διότι τὰ ἀνωτάτω σώματα ἡμῶν θριξὶ περιπέφρακται ; Eust.~1570.12 (ad Od.~7.87) Ὁ δὲ θριγκὸς ἤγουν ἡ τῶν τοίχων κόμη ἢ στεφάνη ἢ στεφάνωσις, παρὰ τὴν τρίχα γίνεται, τριχός ; Eur.~\textit{IT}~47 φεύγειν δὲ κἄξω στᾶσα θριγκὸν εἰσιδεῖν ; Eur.~\textit{Tro.}~489 τὸ λοίσθιον δέ, θριγκὸς ἀθλίων κακῶν ; Aesch.~\textit{Ag.}~1283 κάτεισιν, ἄτας τάσδε θριγκώσων φίλοις ; \textit{Od.}~17.267 τοίχῳ καὶ θριγκοῖσι, θύραι δ᾽ εὐερκέες εἰσὶ ; Pl.~\textit{R.}~7.534e ὥσπερ θριγκὸς τοῖς μαθήμασιν ἡ διαλεκτικὴ ἡμῖν ἐπάνω κεῖσθαι ; Arist.~\textit{Ph.}~7.3,~246a17–20 ἄτοπον γὰρ εἰ ὁ θριγκὸς καὶ ὁ κέραμος ἀλλοίωσις, ἢ εἰ θριγκουμένη καὶ κεραμουμένη ἀλλοιοῦται ἀλλὰ μὴ τελειοῦται ἡ οἰκία ; Luc.~\textit{Gall.}~22 λῃστὴν μὴ ὑφέληται τὸ χρυσίον ὑπερβὰς τὸ θριγκίον ἢ διορύξας τὸν τοῖχον ; Hsch.~γ~267 γεῖσα· τὸ ἐν ταῖς οἰκίαις, ἢ τειχῶν ἐξέχον κυμάτιον (Eur. Phoen. 1158) ; Poll.~7.120 περιοικοδομεῖν θρίγκους, περιβόλους· καὶ θριγκῶσαι τὸ ῥῆμα}
\lp{58–60 θριγκοῖς … κατὰ τοῦ θριγκοῦ]}{Eur.~\textit{Ion}~1320–1321 ἐπίσχες, ὦ παῖ· τρίποδα γὰρ χρηστήριον λιποῦσα θριγκοῦ τοῦδ᾽ ὑπερβάλλω πόδα ; Eur.~\textit{Or.}~1569 ἢ τῷδε θριγκῷ κρᾶτα συνθραύσω σέθεν ; Phot. s.v.~θριγκός (I, p.~283 Naber) θριγκός· τὸ περίφραγμα, στεφάνη, μικρὸν τειχίον, περίβολον ; EM~455.52–57 θριγκὸς οὖν τὸ ἀνώτατον τῆς τοῦ τοίχου οἰκοδομήσεως, καὶ ὁ ὑπὸ τὸν κλινόποδα τοῦ τοίχου τόπος· καὶ ἡ στεφάνη, καὶ τὸ περιχαράκωμα}
\lp{59–62 ὃς ἕτοιμος … χοανεῦσαι]}{Ar.~\textit{Ach.}~274–275 μέσην λαβόντ᾽ ἄραντα καταβαλόντα καταγιγαρτίσ᾽ ; Ar.~\textit{Eq.}~364 ἐγὼ δὲ κινήσω γέ σου τὸν πρωκτὸν ἀντὶ φύσκης ; Ar.~\textit{Av.}~1253–1256 σὺ δ᾽ εἴ με λυπήσεις τι, τῆς διακόνου πρώτης ἀνατείνας τὼ σκέλει διαμηριῶ τὴν Ἶριν αὐτήν, ὥστε θαυμάζειν ὅπως οὕτω γέρων ὢν στύομαι τριέμβολον ; Ar.~\textit{Thesm.}~157–158 ὅταν σατύρους τοίνυν ποιῇς, καλεῖν ἐμέ, ἵνα συμποιῶ σοὔπισθεν ἐστυκὼς ἐγώ ; Catull.~16.1 Pedicabo ego vos et irrumabo ; Ar.~\textit{Nub.}~193–197 τί δῆθ᾽ ὁ πρωκτὸς ἐς τὸν οὐρανὸν βλέπει; … μήπω γε μήπω γ᾽· ἀλλ᾽ ἐπιμεινάντων, ἵνα αὐτοῖσι κοινώσω τι πραγμάτιον ἐμόν ; Ar.~\textit{Pax}~711 ὦ δέσποθ᾽ Ἑρμῆ τῆς Ὀπώρας κατελάσας ; Ar.~\textit{Thesm.}~35 καὶ μὴν βεβίνηκας σύ γ᾽, ἀλλ᾽ οὐκ οἶσθ᾽ ἴσως ; Ar.~\textit{Thesm.}~200–201 ὦ κατάπυγον εὐρύπρωκτος εἶ οὐ τοῖς λόγοισιν ἀλλὰ τοῖς παθήμασιν ; Suid.~κ~498 Καταγιγαρτίσαι: συνουσιάσαι. ἢ ἀντὶ τοῦ κατὰ τῶν γεωργικῶν γιγάρτων βαλεῖν καὶ καταμηρίσαι}
\lp{59 ὃς ἕτοιμος]}{Ar.~\textit{Nub.}~1225–1226 ἵππον; οὐκ ἀκούετε; ὃν πάντες ὑμεῖς ἴστε μισοῦνθ᾽ ἱππικήν ; Soph.~\textit{Phil.}~90 ἀλλ᾽ εἴμ᾽ ἑτοῖμος πρὸς βίαν τὸν ἄνδρ᾽ ἄγειν ; Plb.~12.15.2 κατὰ τὴν πρώτην ἡλικίαν κοινὸν πόρνον, ἕτοιμον τοῖς ἀκρατεστάτοις, κολοιόν, τριόρχην, πάντων τῶν βουλομένων τοῖς ὄπισθεν ἔμπροσθεν γεγονότα ; Hsch.~ε~6630 ἕτοιμος· εὔτροπος. ηὐτρεπισμένος (Σ 96). ταχύς ; Ar.~\textit{Nub.}~806 ἕτοιμος ὅδ᾽ ἐστὶν ἅπαντα δρᾶν ; Ar.~\textit{Vesp.}~341 ἀλλά μ᾽ εὐωχεῖν ἕτοιμός ἐστ᾽· ἐγὼ δ᾽ οὐ βούλομαι ; Ar.~\textit{Ran.}~860 ἕτοιμός εἰμ᾽ ἔγωγε, κοὐκ ἀναδύομαι ; Eur.~\textit{El.}~796 Αἴγισθ᾽, ἕτοιμοι κοὐκ ἀπαρνούμεσθ᾽, ἄναξ ; Eur.~\textit{Phoen.}~969 θνῄσκειν ἕτοιμος πατρίδος ἐκλυτήριον ; Soph.~\textit{OT}~92 ἕτοιμος εἰπεῖν, εἴτε καὶ στείχειν ἔσω ; Pl.~\textit{Cra.}~384e ἕτοιμος ἔγωγε καὶ μανθάνειν καὶ ἀκούειν}
\lp{60 τοῦ καλλιεποῦς]}{Ar.~\textit{Thesm.}~49 μέλλει γὰρ ὁ καλλιεπὴς Ἀγάθων}
\lp{60 κατὰ τοῦ θριγκοῦ]}{Archil.~fr.~196a.21 West (= S~478.21) [θρ]ιγκοῦ δ᾽ ἔνερθε καὶ πυλέων ὑποφ… ; Pl.~\textit{R.}~3.398a ἀποπέμποιμέν τε εἰς ἄλλην πόλιν μύρον κατὰ τῆς κεφαλῆς καταχέαντες καὶ ἐρίῳ στέψαντες ; Σ Ar.~\textit{Thesm.}~60 ἀντὶ τοῦ ἐπὶ τοῦ θριγκοῦ ; Suid.~κ~516 Καταδακτυλικός: ἀντὶ τοῦ συνουσιαστικὸς κατὰ τοῦ δακτυλίου τοῦ πρωκτοῦ ; Ar.~\textit{Eq.}~1094–1095 εἶτα κατασπένδειν κατὰ τῆς κεφαλῆς ἀρυβάλλῳ ἀμβροσίαν κατὰ σοῦ ; Ar.~\textit{Eq.}~1381 οὔκουν καταδακτυλικὸς σὺ τοῦ λαλητικοῦ ; Ar.~\textit{Thesm.}~58 τίς ἀγροιώτας πελάθει θριγκοῖς ; Apollod. com. fr.~13~Kock οὐδὲν γὰρ αἰσχρόν ἐστιν αὐτὸν ἀποτυχεῖν. — πράττουσι πάντα· τὴν γὰρ αἰσχύνην πάλαι — πᾶσαν ἀπολωλέκασι καθ᾽ ἑτέρας θύρας ; Aeschin.~1.80 εἰ γὰρ μνησθείη τειχῶν ἐπισκευῆς ἢ πύργου, ἢ ὡς ἀπήγετό ποι τις, εὐθὺς ἐβοᾶτε καὶ ἐγελᾶτε, καὶ αὐτοὶ ἐλέγετε τὴν ἐπωνυμίαν τῶν ἔργων ὧν σύνιστε αὐτῷ ; D.~19.197 περιρρήξας τὸν χιτωνίσκον ὁ οἰκέτης ξαίνει κατὰ τοῦ νώτου πολλάς ; Ar.~\textit{Nub.}~177–178 κατὰ τῆς τραπέζης καταπάσας λεπτὴν τέφραν κάμψας ὀβελίσκον ; Ar.~\textit{Vesp.}~713 ὡς νάρκη μου κατὰ τῆς χειρὸς καταχεῖται ; Ar.~\textit{Eccl.}~364 τίς τῶν κατὰ πρωκτὸν δεινός ἐστι τὴν τέχνην}
\lp{61 συγγογγύλας καὶ συστρέψας]}{Ar.~\textit{Lys.}~975 ξυστρέψας καὶ ξυγγογγύλας ; Pl.~\textit{R.}~1.336b οὐκέτι ἡσυχίαν ἦγεν, ἀλλὰ συστρέψας ἑαυτὸν ὥσπερ θηρίον ἧκεν ἐφ᾽ ἡμᾶς ὡς διαρπασόμενος ; Σ Ar.~\textit{Thesm.}~61 συγκάμψας ; Hsch.~γ~764 γογγύλλειν· συστρέφειν ; Hdt.~9.18.1 καὶ οἳ ἀντίοι ἔστησαν πάντῃ συστρέψαντες ἑωυτοὺς καὶ πυκνώσαντες ὡς μάλιστα ; Hp.~\textit{Acut.} (Sp.)~37 ταῦτα λεῖα τρίψας, κατάποτα ξυστρέψας, δίδου ; D.H.~\textit{Lys.}~6 ἡ συστρέφουσα τὰ νοήματα καὶ στρογγύλως ἐκφέρουσα λέξις ; D.H.~\textit{Lys.}~13 τὸ συστρέφειν τε καὶ στρογγυλίζειν τὰ νοήματα ; Hsch.~ε~937 εἰλυφόων ⟨ἄνεμος⟩· συστρέφων ἄνεμος (Λ 156) ; Hsch.~ε~4705 ἐπιδινήσας· συστρέψας. συσσείσας (Γ 378)}
\lp{61 συγγογγύλας καὶ συστρέψας (θωμοί, ἅλως)]}{Σ Ar.~\textit{Lys.}~973–974 θωμούς: Σωροὺς ξύλων. θωμοὶ δὲ λέγονται οἱ τῶν πυρῶν σωροί. τυφὼς δὲ ἡ ἐξ ἀναθυμιάσεως τῆς γῆς συστροφή, πρὶν ἐκπυρωθῆναι τὸν ἀέρα. — μεγάλῳ τυφῷ: Μεγάλῳ καύσωνι ; \textit{idem} Suid.~θ~429 (καὶ οἱ τῶν πυρῶν σωροί ; πρὶν ἐκπυρωθῆναι ; add. Ἀριστοφάνης· εἴθ᾽ αὐτὴν ὥσπερ τοὺς θωμοὺς μεγάλῳ τύφῳ καὶ πρηστῆρι ξυστρέψας καὶ ξυγγογγυλίσας οἴχοιο φέρων) ; Hsch.~θ~1005 θωμός· σωρὸς σταχύων, ἢ κορμός ; Hsch.~θ~503 θημών· πᾶν τὸ σεσωρευμένον, οἷον ἀχύρων ἢ πυρῶν (ε 368) ; Phot. s.v.~θωμοί (I, p.~287 Naber) θωμοί· οἱ τῶν σπερμάτων σωροί ; \textit{Il.}~5.499–501 ὡς δ᾽ ἄνεμος ἄχνας φορέει ἱερὰς κατ᾽ ἀλωὰς ἀνδρῶν λικμώντων ; \textit{Il.}~13.588–590 ὡς δ᾽ ὅτ᾽ ἀπὸ πλατέος πτυόφιν μεγάλην κατ᾽ ἀλωὴν θρῴσκωσιν κύαμοι μελανόχροες ἢ ἐρέβινθοι πνοιῇ ὕπο λιγυρῇ καὶ λικμητῆρος ἐρωῇ ; \textit{Il.}~20.495–497 ὡς δ᾽ ὅτε τις ζεύξῃ βόας ἄρσενας εὐρυμετώπους τριβέμεναι κρῖ λευκὸν ἐϋκτιμένῃ ἐν ἀλωῇ, ῥίμφά τε λέπτ᾽ ἐγένοντο βοῶν ὑπὸ πόσσ᾽ ἐριμύκων ; Hes.~\textit{Op.}~597–599 δμωσὶ δ᾽ ἐποτρύνειν Δημήτερος ἱερὸν ἀκτὴν δινέμεν, εὖτ᾽ ἂν πρῶτα φανῇ σθένος Ὠαρίωνος, χώρῳ ἐν εὐαέι καὶ ἐυτροχάλῳ ἐν ἀλωῇ ; \textit{AP}~6.53 εὐξαμένῳ γάρ οἱ ἦλθε βοαθόος, ὄφρα τάχιστα λικμήσῃ πεπόνων καρπὸν ἀπ᾽ ἀσταχύων ; \textit{AP}~6.95 καὶ παλινουροφόρον, χεῖρα θέρευς τρίνακα ; \textit{AP}~6.104 κέντρα τ᾽ ὀπισθονυγῆ, καὶ βουστρόφα δεσμὰ τενόντων, καὶ τρίνακας ξυλίνας, χεῖρας ἀρουροπονων ; \textit{AP}~9.301 τίπτε τὸν ὀγκητὴν βραδύπουν ὄνον ἄμμιγ᾽ ἐν ἵπποις γυρὸν ἀλωειναῖς ἐξελάατε δρόμον ; X.~\textit{Oec.}~18.3–4 οὐκοῦν, ἔφη, τοῦτο μὲν οἶσθα, ὅτι ὑποζυγίῳ ἁλοῶσι τὸν σῖτον ; Paus.~1.38.6 ἐνταῦθα ἅλως καλουμένη Τριπτολέμου καὶ βωμὸς δείκνυται ; Hsch.~θ~501 θημωνιαί· οἱ σωροὶ τῶν δραγμάτων}
\lp{61 συστρέψας]}{Hsch.~σ~2012 στρόβιλος· συστροφὴ ἀελλώδης. καὶ ἡ πίτυς, καὶ ὁ καρπὸς αὐτῆς. καὶ εἶδός τι ὀρχήσεως ; Poll.~5.168 Συστρέψαι τὸ σῶμα, ἀναγαγεῖν, συνελκύσαι, συναγαγεῖν, συναθροῖσαι, συστεῖλαι, συσπειρᾶσαι, συσπειράσασθαι, συγκάμψαι ; Phot.~σ~628 στροβῆσαι· συστρέψαι· καὶ στρόβ†η†λος· ἡ συστροφὴ καὶ ὁ λαῖλαψ· „στρόβ†η†λος ἐν ᾧ συστρέφεται“ (Men. fr. 420, 4 K.- A.) ; Hsch.~ε~3310 ἐνστροβιλίσας· συστρέψας ; Hsch.~σ~2016 στροβῆσαι· συστρέψαι. ἐκφοβῆσαι ; Suid.~ε~863 Ἐλίξας: συστρέψας}
\lp{62 τουτὶ τὸ πέος]}{Ar.~\textit{Nub.}~734 οὐδέν γε πλὴν ἢ τὸ πέος ἐν τῇ δεξιᾷ ; Ar.~\textit{Lys.}~979 περὶ τὴν ψωλὴν περιβαίη ; Ar.~\textit{Lys.}~1012 πρέσβεις ἑλέσθαι τὸ πέος ἐπιδείξας τοδί ; Ar.~\textit{Thesm.}~142 καὶ ποῦ πέος; ποῦ χλαῖνα; ποῦ Λακωνικαί ; Ar.~\textit{Thesm.}~643 ἀνίστασ᾽ ὀρθός. ποῖ τὸ πέος ὠθεῖς κάτω ; Suid.~π~987 Πέος: τὸ σπήλαιον. καὶ Πέος δὲ τὸ αἰδοῖον. Ἀριστοφάνης Νεφέλαις· πλὴν ἢ τὸ πέος ἔχοντα ἐν τῇ δεξιᾷ ; Ar.~\textit{Ach.}~158 τίς τῶν Ὀδομάντων τὸ πέος ἀποτεθρίακεν ; Ar.~\textit{Eq.}~964 ψωλὸν γενέσθαι δεῖ σε μέχρι τοῦ μυρρίνου ; \textit{idem} App. Prov.~5.41 (μεχρὶ τοῦ μυῤῥίνου ; add. μέμνηται ταύτης Ἀριστοφάνης ἐν Ἱππεῦσιν) ; Ar.~\textit{Lys.}~415 νεανίαν καὶ πέος ἔχοντ᾽ οὐ παιδικόν ; Ar.~\textit{Lys.}~928 ἀλλ᾽ ἢ τὸ πέος τόδ᾽ Ἡρακλῆς ξενίζεται ; Ar.~\textit{Thesm.}~254 νὴ τὴν Ἀφροδίτην ἡδύ γ᾽ ὄζει ποσθίου ; Ar.~\textit{Thesm.}~648 τὸ πέος διέλκεις πυκνότερον Κορινθίων ; Hsch.~τ~1216 τουτί· τοῦτο}
\lp{62 χοανεῦσαι]}{\textit{Test.} Moer. χ~21 (p.~213,~11~Bekker) χοανεῦσαι Ἀττικοί, χωνεῦσαι Ἕλληνες ; Hsch.~χ~575 χοανεῦσαι· χωνεῦσαι (Ar. Thesm. 62) ; Ar.~\textit{Lys.}~408–413 ὦ χρυσοχόε τὸν ὅρμον ὃν ἐπεσκεύασας … ἡ βάλανος ἐκπέπτωκεν ἐκ τοῦ τρήματος … σὺ δ᾽ ἢν σχολάσῃς, πάσῃ τέχνῃ πρὸς ἑσπέραν / ἐλθὼν ἐκείνῃ τὴν βάλανον ἐνάρμοσον ; Ar.~\textit{Thesm.}~57 καὶ χοανεύει. — καὶ λαικάζει ; Diph.~fr.~50~Kock (Ath.~10.451b–c) φάσκειν τε τὸν χαλκέα πολὺ κρείττω φέρειν ἰσχύν· ἐπεὶ τοῦτον κατεργαζόμενον καὶ τὸν σίδαρον τὸν σφοδρὸν κάμπτειν, μαλάσσειν, ὅ τι ἂν χρήζῃ ποεῖν. τὰν δὲ τρίταν ἀποφῆναι πέος ἰσχυρότατον πάντων, διδάσκειν δ᾽ ὅτι καὶ τὸν χαλκέα στένοντα πυγίζουσι τούτῳ ; Σ Ar.~\textit{Thesm.}~62 ὡς εἰς χόανον τὸν πρωκτὸν ἐμβαλεῖν ; Ar.~\textit{Thesm.}~18–19 ἀκοῇ δὲ χοάνην ὦτα διετετρήνατο … διὰ τὴν χοάνην οὖν μήτ᾽ ἀκούω μήθ᾽ ὁρῶ ; Pherecr.~fr.~108.30–31~Kock (Ath.~6.268e) πλήρεις κύλικας οἴνου μέλανος ἀνθοσμίου / ἤντλουν διὰ χώνης τοῖσι βουλομένοις πιεῖν ; \textit{idem} Poll.~10.75 (Φερεκράτης ἐν Πετάλῃ … ἤντλουν διὰ χώνης tantum) ; Soph.~fr.~35~R. (=~32 N.) ἀσπὶς μὲν ἡμὴ λίγδος ὣς πυκνομματεῖ ; D.~22.70 συγχωνεύειν ἔπεισεν ; Hsch.~χ~529 †χλοάνοις†· τοῖς κοιλώμασιν, ἐν οἷς ἀναχεῖται τὰ χωνευόμενα ; Poll.~7.106–107 τὰ δὲ ἐργαλεῖα τούτων φῦσαι ἀκροφύσια, χοάναι, πυράγραι, ἄκμονες … τὸ δὲ ἔργον ὑγραίνειν, λύειν, συγχωνεύειν, ἐλαύνειν ἐξελαύνειν, τυποῦν ; Moer. χ~1 (p.~212,~27~Bekker) χοάνην Ἀττικοί, χώνην Ἕλληνες ; Suid.~χ~363 Χοάνης: χωνίου. καὶ Χόανα, τὰ μέταλλα ; Nic.~\textit{Alex.}~52 ἔντοσθεν χοάνοιο διχῆ πυρὸς ἤλασε λιγνύς ; Demetr.~\textit{Eloc.}~281 ὁ τὰς Νίκας τὰς χρυσᾶς χωνεύειν κελεύων ; Din.~1.69 καταχωνεύειν τὸν ἴδιον κόσμον τῶν γυναικῶν ; Lycurg.~\textit{Leoc.}~117 τὴν εἰκόνα αὐτοῦ ἐξ ἀκροπόλεως καθελόντες καὶ συγχωνεύσαντες καὶ ποιήσαντες στήλην ; Hsch.~α~8496 αὐτοχόωνον· αὐτοχώνευτον. ἀπ᾽ αὐτῆς τῆς χωνείας τὸ τέλειον ἔχοντα ; Hsch.~β~368 βαύνη· κάμινος ἢ χωνευτήριον ; Hsch.~ε~4248 ἐπάσκιον· ἡ χώνη. Σικελοί ; Hsch.~κ~393 κάλαθος· … καὶ ἀγγεῖον, ἐν ᾧ χωνεύουσι σίδηρον ; Poll.~10.147 χαλκέως σκεύη ἄκμων ἀκμοθέτης, ῥαιστήρ, πυράγρα, φῦσαι φυσητήρ ἀκροφύσιον, χοάναι, ἀκόναι θηγάναι ; EM~813.13~Sylb. Χόανος] Φῦσαι δ᾽ ἐν χοάνοισιν εἴκοσι πᾶσαι ἐφύσων}
\lp{63 νέος γ᾽ ὢν ἦσθ᾽ ὑβριστής]}{Plaut.~\textit{Amph.}~1031–1032 Prodigum te fuisse oportet olim in adulescentia. … Quia senecta aetate á me mendicas malum ; Soph.~fr.~718~N.² ὕβρις δέ τοι / οὐπώποθ᾽ ἥβης εἰς τὸ σῶφρον ἵκετο, / ἀλλ᾽ ἐν νέοις ἀνθεῖ τε καὶ πάλιν φθίνει ; Pl.~\textit{Euthd.}~273a νεανίσκος τις Παιανιεύς, μάλα καλός τε κἀγαθὸς τὴν φύσιν, ὅσον μὴ ὑβριστὴς [δὲ] διὰ τὸ νέος εἶναι ; Suid.~α~1423 ἐγέννησε δὲ Ἀλυάττην, ὅστις ἕως μὲν νέος ἦν, ὑβριστὴς ἦν καὶ ἀκόλαστος, ἐκβὰς δὲ ἐς ἄνδρα σωφρονέστατος καὶ δικαιότατος ; Ar.~\textit{Nub.}~1417 ἐγὼ δέ γ᾽ ἀντείποιμ᾽ ἂν ὡς δὶς παῖδες οἱ γέροντες ; \textit{idem} Diogenian.~4.18 (Δὶς παῖδες οἱ γέροντες ; add. ἐπὶ τῶν πρὸς τὸ γῆρας εὐηθεστέρων), Tosi n°~1096 (Δὶς παῖδες οἱ γέροντες), Erasm.~\textit{Adag.}~1.5.36 (Δὶς παῖδες οἱ γέροντες, i.e. Bis pueri senes) ; Eur.~\textit{Andr.}~977 ὃ δ᾽ ἦν ὑβριστὴς εἴς τ᾽ ἐμῆς μητρὸς φόνον ; Soph.~\textit{Aj.}~1088 αἴθων ὑβριστής, νῦν δ᾽ ἐγὼ μέγ᾽ αὖ φρονῶ ; X.~\textit{Lac.}~3.2 καταμαθὼν γὰρ τοῖς τηλικούτοις μέγιστον μὲν φρόνημα ἐμφυόμενον, μάλιστα δὲ ὕβριν ἐπιπολάζουσαν ; Arist.~\textit{Rh.}~2.2.6 διὸ οἱ νέοι καὶ οἱ πλούσιοι ὑβρισταί· ὑπερέχειν γὰρ οἴονται ὑβρίζοντες ; Arist.~\textit{Rh.}~2.12.15,~2.13.14 καὶ τὰ ἀδικήματα ἀδικοῦσιν εἰς ὕβριν, οὐ κακουργίαν … καὶ τἀδικήματα ἀδικοῦσιν εἰς κακουργίαν, οὐκ εἰς ὕβριν ; Σ Soph.~\textit{Aj.}~1088 αἴθων ὑβριστής: οὐχ ἵσταται ἐπὶ τῶν αὐτῶν· ἦν τότε ὑβριστὴς οὗτος ἀλλὰ νῦν ἐν μεταβολῇ γέγονεν ; Harp. s.v.~ἔργα νέων παροιμία τίς ἐστιν, ἣν ἀνέγραψε καὶ Ἀριστοφάνης ὁ γραμματικὸς οὕτως ἔχουσαν ἔργα νέων, βουλαὶ δὲ μέσων, εὐχαὶ δὲ γερόντων ; \textit{idem} App. Prov.~4.6 (Νέοις μὲν ἔργα, βουλὰς δὲ γεραιτέροις), Tosi n°~1081 ; Aesch.~\textit{Ag.}~763–766 φιλεῖ δὲ τίκτειν Ὕβρις / μὲν παλαιὰ νεά- / ζουσαν ἐν κακοῖς βροτῶν / ὕβριν ; Soph.~\textit{OC}~395 γέροντα δ᾽ ὀρθοῦν φλαῦρον, ὃς νέος πέσῃ ; Pl.~\textit{Lg.}~884a μέγιστα δὲ δὴ τῶν λοιπῶν αἱ τῶν νέων ἀκολασίαι τε καὶ ὕβρεις ; Str.~14.5.14 ὡς δ᾽ οὐκ ἀπείχοντο ὕβρεως οὐδεμιᾶς … οἱ δὲ πρῶτον μὲν κατετοιχογράφησαν αὐτοῦ τοιαῦτα· ἔργα νέων, βουλαὶ δὲ μέσων, πορδαὶ δὲ γερόντων ; Zen.~2.90 Γέρων ἀλώπηξ οὐχ ἁλίσκεται πάγῃ; ἐπὶ τῶν διὰ χρόνου πλῆθος οὐχ ἁμαρτανόντων}
\lp{63 ὦ γέρον]}{Ar.~\textit{Ach.}~396–397 πῶς ἔνδον εἶτ᾽ οὐκ ἔνδον; — ὀρθῶς ὦ γέρον ; Ar.~\textit{Vesp.}~1299–1303 οὐ γὰρ ὁ γέρων ἀτηρότατον ἄρ᾽ ἦν κακὸν … τούτων ἁπάντων ἦν ὑβριστότατος μακρῷ ; Suid.~κ~2471 Κρόνιππος: ὁ μέγας λῆρος. κατ᾽ ἐπίτασιν λαμβανομένου τοῦ ἵππου. σὺ δ᾽ εἶ κρόνιππος· ἢ ὑβριστής, τρυφητής, πόρνος ; Men. Sent.~146 J. γέρων ἐραστὴς ἐσχάτη κακὴ τύχη ; Ar.~\textit{Nub.}~746 ὦ Σωκρατίδιον φίλτατον. — τί ὦ γέρον ; Ov.~\textit{Am.}~1.9.4 Turpe senex miles, turpe senilis amor ; Diogenian.~1.86 Ἀνδρὸς γέροντος μήποτ᾽ ἐς πυγὴν ὁρᾶν: ἐπὶ τῶν πρὸς ἔνια μὴ χρησίμων ; Diogenian.~2.72 Ἀεὶ γέροντι νέαν ἐπιβάλλειν κούρην: τοῖς πρεσβυτέροις νέαν συνευνάζειν}
\lp{63 ἦ που]}{Soph.~\textit{Aj.}~1228–1230 ἦ που τραφεὶς ἂν μητρὸς εὐγενοῦς ἄπο / ὑψήλ᾽ ἐκόμπεις κἀπ᾽ ἄκρων ὡδοιπόρεις ; Ar.~\textit{Vesp.}~725 ἦ που σοφὸς ἦν ὅστις ἔφασκεν ; Ar.~\textit{Lys.}~1089 ἦ που πρὸς ὄρθρον σπασμὸς ὑμᾶς λαμβάνει ; Ar.~\textit{Ran.}~803 ἦ που βαρέως οἶμαι τὸν Αἰσχύλον φέρειν ; Eur.~\textit{Hec.}~775 ὦ τλῆμον· ἦ που χρυσὸν ἠράσθη λαβεῖν ; Eur.~\textit{Med.}~1308 τί δ᾽ ἔστιν; ἦ που κἄμ᾽ ἀποκτεῖναι θέλει ; Lys.~13.57 καίτοι εἰ ἐκεῖνος ἀπέθανεν, ἦ που Ἀγόρατός γε δικαίως ἀποθανεῖται ; Hsch.~η~705 ἦ που· ὄντως που (Γ 43) ; Ar.~\textit{Plut.}~832 ἦ πού σε ταχέως ἐπέλιπεν τὰ χρήματα ; Eur.~\textit{Hel.}~1465 ἦ που κόρας ἂν ποταμοῦ ; Eur.~\textit{Phoen.}~378 ἦ που στένουσι τλήμονες φυγὰς ἐμάς ; Aesch.~\textit{PV}~521 ἦ πού τι σεμνόν ἐστιν ὃ ξυναμπέχεις ; Σ Eur.~\textit{Phoen.}~1072 ἦ που : ἄρα που. M. ἆρα ποῦ. B}
\lp{63 ὑβριστής]}{Pl.~\textit{Smp.}~175e ὑβριστὴς εἶ, ἔφη, ὦ Σώκρατες, ὁ Ἀγάθων ; Ar.~\textit{Ach.}~479 ἁνὴρ ὑβρίζει· κλῇε πηκτὰ δωμάτων ; Ar.~\textit{Nub.}~1068–1069 κᾆτ᾽ ἀπολιποῦσά γ᾽ αὐτὸν ᾤχετ᾽· οὐ γὰρ ἦν ὑβριστὴς / οὐδ᾽ ἡδὺς ἐν τοῖς στρώμασιν τὴν νύκτα παννυχίζειν ; Eur.~\textit{Ba.}~743 ταῦροι δ᾽ ὑβρισταὶ κἀς κέρας θυμούμενοι ; X.~\textit{Cyr.}~7.5.62 οἵ τε ὑβρισταὶ ἵπποι ἐκτεμνόμενοι τοῦ μὲν δάκνειν καὶ ὑβρίζειν ἀποπαύονται ; Aeschin.~1.15 ἐάν τις ὑβρίζῃ εἰς παῖδα (ὑβρίζει δὲ δή που ὁ μισθούμενος) ἢ ἄνδρα ἢ γυναῖκα ; Hsch.~α~4253 ἀνάγωγος· ὑβριστής. ὁ ⟨μὴ⟩ τῆς δεούσης ἀγωγῆς τετυχηκώς ; Hsch.~ε~4916 ἐπίκωμος· ὁ ἐπᾴδων τῇ φίλῃ. ἢ ἀπὸ οἴνου ᾄδων. ὑβριστής. συγχαίρων ; Hsch.~μ~316 μάρπτις· ὑβριστής (Aesch. Suppl. 826) ; Hsch.~υ~35 ὑβριστής· ὑπερήφανος, ἀγνώμων ; Phot.~υ~11 ὕβρις· ἡ μετὰ προπηλακισμοῦ καὶ ἐπηρείας {καὶ} αἰκία· αἰκία δὲ πληγαὶ μόνον ; Greg. Cypr.~3.86 Ὕβρις ὕβριν ἔτικτε καὶ ψόγος ψόγον ; Ar.~\textit{Ach.}~1117 οἴμ᾽ ὡς ὑβρίζεις ; Ar.~\textit{Pax}~1264 ὑβριζόμεθα. χωρῶμεν ὦ τᾶν ἐκποδών ; Ar.~\textit{Thesm.}~465 δεῖ δὲ ταύτης τῆς ὕβρεως ἡμῖν / τὸν ἄνδρα ; Ar.~\textit{Thesm.}~535 ταύτην ἐῶσαι τὴν φθόρον τοιαῦτα περιυβρίζειν ; Ar.~\textit{Ran.}~21 εἶτ᾽ οὐχ ὕβρις ταῦτ᾽ ἐστὶ καὶ πολλὴ τρυφή ; Ar.~\textit{Plut.}~899 τούτους ὑβρίζειν εἰς ἔμ᾽; οἴμ᾽ ὡς ἄχθομαι ; Ar.~\textit{Plut.}~1044 τάλαιν᾽ ἐγὼ τῆς ὕβρεος ἧς ὑβρίζομαι ; Ar.~\textit{Plut.}~1074 εἶναί σ᾽ ὑβριστήν φησι ; \textit{Od.}~6.120 ἦ ῥ᾽ οἵ γ᾽ ὑβρισταί τε καὶ ἄγριοι οὐδὲ δίκαιοι ; Th.~8.74.3 καὶ ὅτι αὐτῶν καὶ γυναῖκες καὶ παῖδες ὑβρίζονται ; X.~\textit{Mem.}~1.2.12 Ἀλκιβιάδης δὲ αὖ τῶν ἐν τῇ δημοκρατίᾳ πάντων ἀκρατέστατός τε καὶ ὑβριστότατος ; Arist.~\textit{EN}~1148b29–30 ἡ τῶν ἀφροδισίων τοῖς ἄρρεσιν· τοῖς μὲν γὰρ φύσει τοῖς δ᾽ ἐξ ἔθους συμβαίνουσιν, οἷον τοῖς ὑβριζομένοις ἐκ παίδων ; Lys.~12.98 οἱ δὲ παῖδες ὑμῶν, ὅσοι μὲν ἐνθάδε ἦσαν, ὑπὸ τούτων ἂν ὑβρίζοντο ; Hsch.~λ~1487 λωβητήρ· ὑβριστής (Λ 385) ; Hsch.~μ~272 μαργοῖς· μαινομένοις. ὑβρισταῖς (Plat. leg. 792e) ; Hsch.~υ~36 ὑβριστοδίκαι· οἱ μὴ θέλοντες εἰσάγειν τὰς δίκας παρὰ Ἀττικοῖς ; Suid.~π~736 Παροίνιος: ὁ μέθυσος καὶ ὑβριστής ; Diogenian.~5.79 Κολοφωνία ὕβρις: ἐπὶ τῶν πλουσίων καὶ ὑβριστῶν. Τοιοῦτοι γὰρ οἱ Κολοφώνιοι}
\lp{64 ὦ δαιμόνιε]}{Ar.~\textit{Nub.}~38 ἔασον ὦ δαιμόνιε καταδαρθεῖν τί με ; Hsch.~δ~74 δαιμόνιε· μακάριε. κακόδαιμον ; \textit{idem} Suid.~δ~115 (om. κακόδαιμον) ; Ar.~\textit{Eq.}~860 ὦ δαιμόνιε μὴ τοῦ λέγοντος ἴσθι ; Ar.~\textit{Vesp.}~962 ἄκουσον ὦ δαιμόνιέ μου τῶν μαρτύρων ; Ar.~\textit{Av.}~1436 ὦ δαιμόνιε μὴ νουθέτει μ᾽ ἀλλὰ πτέρου ; Ar.~\textit{Lys.}~945 ἀγαθόν· ἔα αὔτ᾽ ὦ δαιμονία. — ληρεῖς ἔχων ; Ar.~\textit{Ran.}~44 ὦ δαιμόνιε πρόσελθε· δέομαι γάρ τί σου ; Ar.~\textit{Eccl.}~564 ὦ δαιμόνι᾽ ἀνδρῶν τὴν γυναῖκ᾽ ἔα λέγειν ; Ar.~\textit{Eccl.}~784 ὦ δαιμόνι᾽ ἀνδρῶν ἔα με τῶν προὔργου τι δρᾶν ; Pherecr.~fr.~80.1~Kock (Ath.~3.75b) ὦ δαιμόνιε, πύρεττε μηδὲν φροντίσας ; Ar.~\textit{Nub.}~816 ὦ δαιμόνιε, τί χρῆμα πάσχεις ὦ πάτερ ; Ar.~\textit{Nub.}~1138 ὦ δαιμόνιε τὸ μέν τι νυνὶ μὴ λάβῃς ; Ar.~\textit{Av.}~961 ὦ δαιμόνιε τὰ θεῖα μὴ φαύλως φέρε ; Ar.~\textit{Av.}~1638 ὦ δαιμόνι᾽ ἀνθρώπων Πόσειδον ποῖ φέρει ; Mnesim.~fr.~3.4~Kock (Ath.~8.359c) πῶς ἔτι μετριώτερ᾽, ὦ δαιμόνι᾽ ; \textit{Il.}~1.561 δαιμονίη αἰεὶ μὲν ὀΐεαι οὐδέ σε λήθω ; \textit{Il.}~2.190 δαιμόνι᾽ οὔ σε ἔοικε κακὸν ὣς δειδίσσεσθαι ; \textit{Od.}~10.472 δαιμόνι᾽, ἤδη νῦν μιμνήσκεο πατρίδος αἴης ; Hdt.~4.126.1 δαιμόνιε ἀνδρῶν, τί φεύγεις αἰεί ; Hdt.~7.48.1 δαιμόνιε ἀνδρῶν, κοῖα ταῦτα λέγεις ; Hdt.~8.84.2 ὦ δαιμόνιοι, μέχρι κόσου ἔτι πρύμνην ἀνακρούεσθε ; Pl.~\textit{Grg.}~489d ἀλλ᾽ ἐγὼ μέν, ὦ δαιμόνιε, καὶ αὐτὸς πάλαι τοπάζω ; Pl.~\textit{Phdr.}~235c πλῆρές πως, ὦ δαιμόνιε, τὸ στῆθος ἔχων}
\lp{64 τοῦτον μὲν ἔα χαίρειν, σὺ δέ]}{Ar.~\textit{Ach.}~200 χαίρειν κελεύων πολλὰ τοὺς Ἀχαρνέας ; Ar.~\textit{Nub.}~932 δεῦρ᾽ ἴθι, τοῦτον δ᾽ ἔα μαίνεσθαι ; Eur.~fr.~23 N. ἀλλ᾽ ἢ τὸ γῆρας τὴν Κύπριν χαίρειν ἐᾷ, / ἥ τ᾽ Ἀφροδίτη τοῖς γέρουσιν ἄχθεται ; Soph.~\textit{OT}~1069–1070 ἄξει τις ἐλθὼν δεῦρο τὸν βοτῆρά μοι; / ταύτην δ᾽ ἐᾶτε πλουσίῳ χαίρειν γένει ; Pl.~\textit{Euthd.}~307b ἀλλ᾽ ἐάσας χαίρειν τοὺς ἐπιτηδεύοντας φιλοσοφίαν ; D.~45.85 ἐγὼ δὲ τούτῳ μὲν χαίρειν λέγω, οὓς δ᾽ ὁ πατήρ μοι παρέδωκε βοηθοὺς καὶ φίλους, εἰς τούτους ἥκω ; Diogenian.~7.9 Οὐκ ἔστ᾽ ἐμὸν τὸ πρᾶγμα, πολλὰ χαιρέτω: ἐπὶ τῶν ἀπραγμόνων ; Ar.~\textit{Plut.}~1187 (τὸν οὖν Δία τὸν σωτῆρα) χαίρειν ἐάσας ἐνθάδ᾽ αὐτοῦ καταμένειν ; Stratt.~fr.~41.2~Kock (Σ Ar.~\textit{Vesp.}~1346) ἦ μήποτ᾽, ὦ παῖ Ζηνός, ἐς ταὐτὸν μόλῃς, / ἀλλὰ παραδοὺς τοῖς Λεσβίοις χαίρειν ἔα ; Soph.~\textit{El.}~448 οὐχ ἔστιν. ἀλλὰ ταῦτα μὲν μέθες· σὺ δὲ ; Pl.~\textit{Phd.}~63e καὶ ὁ Σωκράτης, ἔα, ἔφη, χαίρειν αὐτόν· ἀλλὰ μόνον τὸ ἑαυτοῦ παρασκευαζέτω ; Pl.~\textit{Smp.}~176e τὴν μὲν ἄρτι εἰσελθοῦσαν αὐλητρίδα χαίρειν ἐᾶν … ἡμᾶς δὲ διὰ λόγων ἀλλήλοις συνεῖναι τὸ τήμερον ; Pl.~\textit{Sph.}~239c σὲ μὲν καὶ ἐμὲ χαίρειν ἐῶμεν ; Luc.~\textit{DDeor.}~10.3 ἐκεῖνα μὲν πάντα χαίρειν ἔα καὶ ἐπιλάθου αὐτῶν … σὺ δὲ ; Hdt.~4.112.1 ἐπεὶ δὲ ἔμαθον αὐτοὺς αἱ Ἀμαζόνες ἐπ᾽ οὐδεμιῇ δηλήσι ἀπιγμένους, ἔων χαίρειν ; Pl.~\textit{Cri.}~45a εἰ γάρ τι τοιοῦτον φοβῇ, ἔασον αὐτὸ χαίρειν ; Pl.~\textit{Ep.}~2.312b εἰ μὲν ὅλως φιλοσοφίας καταπεφρόνηκας, ἐᾶν χαίρειν ; Pl.~\textit{Prt.}~347b περὶ μὲν ᾀσμάτων τε καὶ ἐπῶν ἐάσωμεν, περὶ δὲ ὧν τὸ πρῶτον ἐγώ σε ἠρώτησα ; Plu.~\textit{Alc.}~12.3 τὸν δὲ Διομήδη χαίρειν εἴασε χαλεπῶς φέροντα ; Luc.~\textit{DMeretr.}~1.2 καὶ σὺ ἄλλον, ὦ Γλυκέριον, τρυγήσεις, τοῦτον δὲ χαίρειν ἔα}
\lp{64 ἔα χαίρειν]}{Eur.~\textit{El.}~400 βροτῶν δὲ μαντικὴν χαίρειν ἐῶ ; Hsch.~ε~7744 ἐῶν χαίρειν· παραπεμπόμενος ; Hsch.~χ~18 χαίρειν [θ]ἐᾶσαι· καταφρονεῖν (Plat. Soph. 220 a . .) ; Eur.~\textit{Cyc.}~318–319 ἄκρας δ᾽ ἐναλίας ἃς καθίδρυται πατὴρ / χαίρειν κελεύω ; Eur.~\textit{Hipp.}~113 τὴν σὴν δὲ Κύπριν πόλλ᾽ ἐγὼ χαίρειν λέγω ; Eur.~\textit{Hipp.}~1059 φοιτῶντας ὄρνις πόλλ᾽ ἐγὼ χαίρειν λέγω ; Eur.~\textit{Med.}~1044 οὐκ ἂν δυναίμην· χαιρέτω βουλεύματα ; Eur.~\textit{Phoen.}~919 οὐκ ἔκλυον, οὐκ ἤκουσα· χαιρέτω πόλις ; Pl.~\textit{Grg.}~461a εἰ δὲ μή, ἐᾶν χαίρειν ; Pl.~\textit{Phd.}~100d τὰ μὲν ἄλλα χαίρειν ἐῶ ; Pl.~\textit{Phdr.}~230a ὅθεν δὴ χαίρειν ἐάσας ταῦτα ; Pl.~\textit{Sph.}~220a τὸ μὲν τῶν ἀψύχων, ἀνώνυμον ὂν … χαίρειν ἐᾶσαι ; Hsch.~χ~17 χαίρειν· ἐρρῶσθαι. τὸ χαίρειν ταῖς ἐπιστολαῖς προσετίθεσαν. ἔστι δὲ καὶ ἀπαλλασσομένων προσαγόρευσις (Eupol. fr. 331, 2 K.–A.) ; \textit{idem} App. Prov.~5.20 (om. ἐρρῶσθαι ; add. ὡς καὶ Ἀριστοφάνης καὶ Σοφοκλῆς) ; Hsch.~χ~19 χαίρειν φράσαντες· ἀποταξάμενοι, ⟨καταγνόντες⟩ ; \textit{idem} Synag.~χ~3 (vers. ant.), Suid.~χ~168, EM~808.4 ; Hsch.~χ~20 χαιρέτω ληρῶν· ἐρυθριάτω ⟨φλυαρῶν⟩ ; \textit{idem} Synag.~χ~4 (vers. ant.), EM~808.3 ; Synag.~χ~1 (vers. ant.) χαίρειν· οἰμώζειν, στενάζειν ; \textit{idem} Suid.~μ~72 (Μακρὰ χαίρειν), Suid.~χ~161, EM~808.2 ; Suid.~χ~167 Χαίρειν εἰπών: ἀντὶ τοῦ καταλιπών ; Eur.~fr.~1025 N. τὰς θνητῶν δ᾽ ἐγὼ / χαίρειν κελεύω θεῶν ἄτερ προθυμίας ; Aesch.~\textit{Ag.}~572 καὶ πολλὰ χαίρειν ξυμφορὰς καταξιῶ ; Soph.~\textit{El.}~1456 ἦ πολλὰ χαίρειν μ᾽ εἶπας οὐκ εἰωθότως ; Hdt.~9.45.2 τὰ μὲν σφάγια ἐᾶν χαίρειν ; Hsch.~π~2800 πολλὰ χαίρει⟨ν⟩ φράσας· ἀποταξάμενος (Phil. Flacc. 101. leg. Gai. 65) ; \textit{idem} Phot.~π~1024 (add. ἀπογνούς), Suid.~π~1873 (add. ἀπογνούς)}
\lp{64 ἔα]}{Hsch.~ε~4 ἔα· ἄφες (Α 276) ; Hsch.~ε~8 ἔα δή· ἄγε δή (Ar. Thesm. 659) ; Hsch.~ε~22 ἐᾶν· ἀφιέναι (Ο 347)}
\lp{65 Ἀγάθωνά μοι δεῦρ᾽ ἐκκάλεσον]}{Ar.~\textit{Ach.}~395–409 ἐκκάλεσον αὐτόν. — ἀλλ᾽ ἀδύνατον. — ἀλλ᾽ ὅμως· … ἀλλ᾽ ἐκκυκλήθητ᾽. — ἀλλ᾽ ἀδύνατον. — ἀλλ᾽ ὅμως ; Ar.~\textit{Nub.}~218–220a αὐτός. — τίς αὐτός; — Σωκράτης. … ἴθ᾽ οὗτος, ἀναβόησον αὐτόν μοι μέγα. — αὐτὸς μὲν οὖν σὺ κάλεσον· οὐ γάρ μοι σχολή ; Ar.~\textit{Pax}~195 ἴθι νυν κάλεσόν μοι τὸν Δί᾽ ; Eur.~\textit{Ba.}~170 τίς ἐν πύλαισι; Κάδμον ἐκκάλει δόμων ; Ar.~\textit{Vesp.}~271 ᾄδοντας αὐτὸν ἐκκαλεῖν, ἤν τί πως ἀκούσας ; Ar.~\textit{Lys.}~850 πρὸς τῶν θεῶν νυν ἐκκάλεσόν μοι Μυρρίνην ; Ar.~\textit{Eccl.}~34 τήνδ᾽ ἐκκαλέσωμαι θρυγονῶσα τὴν θύραν ; Hsch.~ε~1421 ἐκκαλεῖται· προκαλεῖται ; Zen.~1.89 Ἀνεῳγμέναι Μουσῶν θύραι: ἐπὶ τῶν ἐξ ἑτοίμου λαμβανόντων τὰ κάλλιστα τῶν ἐν παιδείᾳ ; Ar.~\textit{Lys.}~875 οὐ γὰρ δεόμενος οὐδὲν ἐκκαλεῖς ἐμέ ; Ar.~\textit{Thesm.}~289–291 καὶ τοῦ θυγατρίου χοῖρον ἀνδρός μοι τυχεῖν … καὶ ποσθαλίσκον νοῦν ἔχειν μοι καὶ φρένας ; Ar.~\textit{Thesm.}~627 σὺ δ᾽ ἀπόστηθί μοι ; Ar.~\textit{Plut.}~1103 ἀλλ᾽ ἐκκάλει τὸν δεσπότην τρέχων ταχύ ; Soph.~\textit{OT}~597 νῦν οἱ σέθεν χρῄζοντες ἐκκαλοῦσί με ; Soph.~\textit{Phil.}~1264 τί μ᾽ ἐκκαλεῖσθε; τοῦ κεχρημένοι, ξένοι ; Hsch.~ε~4056 ἐπαγγεῖλαι· ἐκκαλέσαι. ἐντυχεῖν}
\lp{65 πάσῃ τέχνῃ]}{Ar.~\textit{Nub.}~885 ἐὰν δὲ μή, τὸν γοῦν ἄδικον πάσῃ τέχνῃ ; Ar.~\textit{Lys.}~412 σὺ δ᾽ ἢν σχολάσῃς, πάσῃ τέχνῃ πρὸς ἑσπέραν ; Ar.~\textit{Ran.}~1235 ἀλλ᾽ ὦγάθ᾽ ἔτι καὶ νῦν ἀπόδος πάσῃ τέχνῃ ; Ar.~\textit{Eccl.}~366 Ἀντισθένη τις καλεσάτω πάσῃ τέχνῃ ; Eur.~\textit{Heracl.}~1002 ἔχθραν πατρῴαν, πάντα κινῆσαι πέτρον ; \textit{idem} Diogenian.~7.42 (Πάντα κινήσω πέτρον ; add. ὁμοία τῇ, Πάντα λίθον κινήσω κτλ.), Hsch.~π~389 (πάντα κινήσας πέτρον· παροιμία), Suid.~π~222 (Πάντα κινήσω πέτρον ; add. ὁμοία τῇ πάντα λίθον κίνει) ; Suid.~δ~1083 Δικαίως ἀδίκως: ἀντὶ τοῦ παντὶ τρόπῳ. Ἀριστοφάνης Πλούτῳ· … καὶ δικαίως καὶ ἀδίκως. τουτέστι πάσῃ τέχνῃ ; Suid.~π~753 Πάσῃ τέχνῃ: ἐκ παντὸς τρόπου. Ἀριστοφάνης Λυσιστράτῃ· οὐδ᾽ ἢν σχολάσῃς, πάσῃ τέχνῃ πρὸς ἑσπέραν ἐλθών, ἐκείνῃ τὴν βάλανον ἐνάρμοσον ; Ar.~\textit{Nub.}~1323 ἀμυνάθετέ μοι τυπτομένῳ πάσῃ τέχνῃ ; Ar.~\textit{Thesm.}~271 πάσαις τέχναις, ἤν μοί τι περιπίπτῃ κακόν ; Ar.~\textit{Eccl.}~534 ἥπερ μεθῆκέ μ᾽ ἐξιέναι πάσῃ τέχνῃ ; Ar.~\textit{Plut.}~233 μεστὴν ποιῆσαι καὶ δικαίως κἀδίκως ; Eur.~\textit{Med.}~278 ἐχθροὶ γὰρ ἐξιᾶσι πάντα δὴ κάλων ; \textit{idem} Ar.~\textit{Eq.}~756 (νῦν δή σε πάντα δεῖ κάλων ἐξιέναι σεαυτοῦ), Zen.~5.62 (Πάντα κάλων σεῖε ; add. παροιμία ἐπὶ τῶν πάσῃ προθυμίᾳ χρωμένων κτλ.) ; Soph.~\textit{Aj.}~752 θεὶς εἶπε κἀπέσκηψε, παντοίᾳ τέχνῃ ; Th.~5.18.4 μήτε τέχνῃ μήτε μηχανῇ μηδεμιᾷ ; D.~24.150 οὔτε τέχνῃ οὔτε μηχανῇ οὐδεμιᾷ ; Lys.~19.11 δέομαι δ᾽ ὑμῶν πάσῃ τέχνῃ καὶ μηχανῇ ; Hsch.~τ~722 τέχνη· ἐπιστήμη (Γ 61). ἢ δόλος (δ 455). καὶ ἀτεχνῶς· ἀδόλως ; Ar.~\textit{Eq.}~592–593 γὰρ τοῖς ἀνδράσι τοῖσδε πάσῃ τέχνῃ πορίσαι σε νίκην ; Ar.~\textit{Lys.}~300 τοῦτο πάσῃ μηχανῇ ; Ar.~\textit{Thesm.}~430 ἢ φαρμάκοισιν ἢ μιᾷ γέ τῳ τέχνῃ ; Hdt.~1.112.1 ἐχρήιζε μηδεμιῇ τέχνῃ ἐκθεῖναί μιν ; X.~\textit{An.}~4.5.16 ἐδεῖτο αὐτῶν πάσῃ τέχνῃ καὶ μηχανῇ μὴ ἀπολείπεσθαι ; X.~\textit{An.}~7.2.8 κελεύει πάσῃ τέχνῃ καὶ μηχανῇ πλεῦσαι ; Lys.~19.53 ὦ ἄνδρες δικασταί, πάσῃ τέχνῃ καὶ μηχανῇ ἐλεήσατε ; Suid.~χ~168 Χαίρειν, ἤν τε βούληθ᾽, ἤν τε μή: ἀντὶ τοῦ παντὶ τρόπῳ ; Zen.~1.93 Ἂν ἡ λεοντῆ μὴ ἐξίκηται, τὴν ἀλωπεκῆν πρόσαψον· … ἂν μὴ κατὰ ῥώμην τὸ προκείμενον ἐξανύοιτο, μηχανῇ καὶ τέχνῃ περαινέσθω τὸ λειπόμενον ; Zen.~5.63 Τὸν δὲ Ἀπόλλωνα ἀποκρίνασθαί φασι, πάντα λίθον κίνει· ὅθεν ἡ παροιμία ἐκράτησεν ; \textit{idem} Hsch.~π~390 (πάντα λίθον κινεῖν· παροιμία, ἧς μέμνηται Ἀρίσταρχος κτλ.), Tosi n°~1286 (Πάντα λίθον κίνει), Erasm.~\textit{Adag.}~1.4.30 (Πάντα λίθον κίνει, i.e. Unumquemque move lapidem)}
